\documentclass{article}

% if you need to pass options to natbib, use, e.g.:
%     \PassOptionsToPackage{numbers, compress}{natbib}
% before loading neurips_2026

% The authors should use one of these tracks.
% Before accepting by the NeurIPS conference, select one of the options below.
% 0. "default" for submission
 \usepackage{neurips_2026}
% the "default" option is equal to the "main" option, which is used for the Main Track with double-blind reviewing.
% 1. "main" option is used for the Main Track
%  \usepackage[main]{neurips_2026}
% 2. "position" option is used for the Position Paper Track
%  \usepackage[position]{neurips_2026}
% 3. "dandb" option is used for the Datasets & Benchmarks Track
 % \usepackage[dandb]{neurips_2026}
% 4. "creativeai" option is used for the Creative AI Track
%  \usepackage[creativeai]{neurips_2026}
% 5. "sglblindworkshop" option is used for the Workshop with single-blind reviewing
 % \usepackage[sglblindworkshop]{neurips_2026}
% 6. "dblblindworkshop" option is used for the Workshop with double-blind reviewing
%  \usepackage[dblblindworkshop]{neurips_2026}

% After being accepted, the authors should add "final" behind the track to compile a camera-ready version.
% 1. Main Track
 % \usepackage[main, final]{neurips_2026}
% 2. Position Paper Track
%  \usepackage[position, final]{neurips_2026}
% 3. Datasets & Benchmarks Track
 % \usepackage[dandb, final]{neurips_2026}
% 4. Creative AI Track
%  \usepackage[creativeai, final]{neurips_2026}
% 5. Workshop with single-blind reviewing
%  \usepackage[sglblindworkshop, final]{neurips_2026}
% 6. Workshop with double-blind reviewing
%  \usepackage[dblblindworkshop, final]{neurips_2026}
% Note. For the workshop paper template, both \title{} and \workshoptitle{} are required, with the former indicating the paper title shown in the title and the latter indicating the workshop title displayed in the footnote.
% For workshops (5., 6.), the authors should add the name of the workshop, "\workshoptitle" command is used to set the workshop title.
% \workshoptitle{WORKSHOP TITLE}

% "preprint" option is used for arXiv or other preprint submissions
 % \usepackage[preprint]{neurips_2026}

% to avoid loading the natbib package, add option nonatbib:
%    \usepackage[nonatbib]{neurips_2026}
\usepackage{graphicx} % Required for inserting images

% --- Theorem environments (Definition / Theorem / Lemma / Claim) ---
\usepackage{amsthm}
\usepackage{amssymb} % provides \mathbb
\usepackage{thmtools}
\usepackage{thm-restate}
\usepackage{booktabs}

%for colored boxes
\usepackage[most]{tcolorbox}
\usepackage{xcolor}

% Define the color used in the paper (a light grey/blue)
\definecolor{messagecolor}{RGB}{245, 245, 245}
\definecolor{bordercolor}{RGB}{180, 180, 180}

%conclusion "bronstein" boxes
% Define the custom box environment
\usepackage[most]{tcolorbox}

% Define the custom box environment
\newtcolorbox{conclusionbox}[1]{
  enhanced,
  arc=4pt,                  
  boxrule=1pt,              
  colframe=cyan!70!black,   
  colback=cyan!4!white,     % <-- CHANGED: Faint cyan/teal background
  coltitle=black,           
  fonttitle=\bfseries,      
  attach boxed title to top left={xshift=15pt, yshift=-6pt}, 
  boxed title style={
    colback=white,          % Keep the title box background solid white for contrast
    colframe=cyan!70!black, 
    arc=3pt,                
    boxrule=1pt             
  },
  title=#1                  
}
%conclusion "bronstein" boxes


\newtcolorbox{messagebox}[1]{
  colback=messagecolor,
  colframe=bordercolor,
  fonttitle=\bfseries,
  coltitle=black,
  title=#1,
  enhanced,
  attach title to upper,
  after title={:\enskip}, % Adds a colon after the title
  sharp corners,
  boxrule=0.5pt,
  left=5pt,
  right=5pt,
  top=5pt,
  bottom=5pt
}


\usepackage{todonotes}
% Numbering: <section>.<number>
%\newcommand{\nd}[1]{\textcolor{blue}{\textit{[Nadav: #1]}}}
\newcommand{\nd}{\todo[color=blue!40]}
\newcommand{\ys}[1]{\textcolor{orange}{\textit{[Yon: #1]}}}
%\newcommand{\bgf}[1]{\textcolor{red}{\textit{[Benjy: #1]}}}
\newcommand{\bgf}{\todo[color=red!40]}
\newcommand{\sout}[1]{\textcolor{pink}{\textit{[Soutrik: #1]}}}

\theoremstyle{definition}
\newtheorem{definition}{Definition}[section]
\newtheorem{remark}[definition]{Remark}
\newtheorem{example}[definition]{Example}
\newtheorem{conclude}[definition]{Conclusion}

\theoremstyle{plain}
\newtheorem{theorem}[definition]{Theorem}
\newtheorem{proposition}[definition]{Proposition}
\newtheorem{lemma}[definition]{Lemma}
\newtheorem{claim}[definition]{Claim}

%\theoremstyle{remark}


\usepackage{amsmath}
\usepackage{mathrsfs}
% Blackboard bold shortcuts:
\newcommand{\bbA}{\mathbb{A}}
\newcommand{\bbB}{\mathbb{B}}
\newcommand{\bbC}{\mathbb{C}}
\newcommand{\D}{\mathcal{D}}
\newcommand{\bbE}{\mathbb{E}}
\newcommand{\bbF}{\mathbb{F}}
\newcommand{\G}{\mathbb{G}}
\newcommand{\bbH}{\mathbb{H}}
\newcommand{\bbI}{\mathbb{I}}
\newcommand{\bbJ}{\mathbb{J}}
\newcommand{\bbK}{\mathbb{K}}
\newcommand{\bbL}{\mathbb{L}}
\newcommand{\bbM}{\mathbb{M}}
\newcommand{\N}{\mathcal{N}}
\newcommand{\bbN}{\mathbb{N}}
\newcommand{\bbO}{\mathbb{O}}
\newcommand{\bbP}{\mathbb{P}}
\newcommand{\bbQ}{\mathbb{Q}}
\newcommand{\R}{\mathbb{R}}
\newcommand{\RR}{\mathbb{R}}
\newcommand{\bbS}{\mathbb{S}}
\newcommand{\bbT}{\mathbb{T}}
\newcommand{\bbU}{\mathbb{U}}
\newcommand{\bbV}{\mathbb{V}}
\newcommand{\bbW}{\mathbb{W}}
\newcommand{\X}{\mathcal{X}}
\newcommand{\Y}{\mathcal{Y}}
\newcommand{\bbZ}{\mathbb{Z}}
%\usepackage[margin=1in]{geometry}
%\newtheorem{proposition}{Proposition}

\newcommand{\csort}{c_{\mathrm{lex}}}
\newcommand{\Lemp}{\mathcal{L}_{\mathrm{emp}}}

\newcommand{\twopartdef}[4]
{
	\left\{
		\begin{array}{ll}
			#1 & \mbox{if } #2 \\
			#3 & \mbox{if } #4
		\end{array}
	\right.
}

\newcommand{\threepartdef}[6]
{
	\left\{
		\begin{array}{lll}
			#1 & \mbox{if } #2 \\
			#3 & \mbox{if } #4 \\
			#5 & \mbox{if } #6
		\end{array}
	\right.
}


\newcommand{\Sym}{\mathrm{Sym}}
\newcommand{\rank}{\mathrm{rank}}
\newcommand{\diag}{\mathrm{diag}}
\newcommand{\tr}{\mathrm{tr}}
\newcommand{\eps}{\varepsilon}
\usepackage[utf8]{inputenc} % allow utf-8 input
\usepackage[T1]{fontenc}    % use 8-bit T1 fonts
\usepackage{hyperref}       % hyperlinks
\usepackage{url}            % simple URL typesetting
\usepackage{booktabs}       % professional-quality tables
\usepackage{amsfonts}       % blackboard math symbols
\usepackage{nicefrac}       % compact symbols for 1/2, etc.
\usepackage{microtype}      % microtypography
\usepackage{xcolor}         % colors

% Note. For the workshop paper template, both \title{} and \workshoptitle{} are required, with the former indicating the paper title shown in the title and the latter indicating the workshop title displayed in the footnote. 
\title{When and How to Canonize: a Generalization Perspective}


% The \author macro works with any number of authors. There are two commands
% used to separate the names and addresses of multiple authors: \And and \AND.
%
% Using \And between authors leaves it to LaTeX to determine where to break the
% lines. Using \AND forces a line break at that point. So, if LaTeX puts 3 of 4
% authors names on the first line, and the last on the second line, try using
% \AND instead of \And before the third author name.


\author{%
  David S.~Hippocampus\thanks{Use footnote for providing further information
    about author (webpage, alternative address)---\emph{not} for acknowledging
    funding agencies.} \\
  Department of Computer Science\\
  Cranberry-Lemon University\\
  Pittsburgh, PA 15213 \\
  \texttt{hippo@cs.cranberry-lemon.edu} \\
  % examples of more authors
  % \And
  % Coauthor \\
  % Affiliation \\
  % Address \\
  % \texttt{email} \\
  % \AND
  % Coauthor \\
  % Affiliation \\
  % Address \\
  % \texttt{email} \\
  % \And
  % Coauthor \\
  % Affiliation \\
  % Address \\
  % \texttt{email} \\
  % \And
  % Coauthor \\
  % Affiliation \\
  % Address \\
  % \texttt{email} \\
}


\begin{document}


\maketitle


\begin{abstract}

While equivariant architectures are standard for processing symmetric data, there is growing interest in achieving equivariance by applying group averaging or canonization to non-equivariant backbones. However, the theoretical generalization properties of these alternative strategies remain poorly understood. We introduce a theoretical framework to analyze the generalization error of these methods by bounding their covering numbers. We establish a rigorous generalization hierarchy: the error bounds of canonized models are at best equal to the error bounds of structurally equivariant and group-averaged models, and at worst equal to the bounds of  non-equivariant baselines. Furthermore, we show that there exist "optimal" canonizations which attain the optimal error bounds, and "poor" canonizations which attain the non-equivariant error bounds, and that this depends on the regularity of the canonization. Finally, applying this framework to permutation groups in point cloud processing, we rigorously prove that the covering number of lexicographical sorting grows exponentially with point cloud dimension, whereas Hilbert curve canonization guarantees polynomial growth. This provides the first formal theoretical justification for the empirical success of Hilbert curve serialization in state-of-the-art point cloud architectures. We conclude with experiments that support our theoretical claims. %Ultimately, this work provides the first formal theoretical justification for the empirical success of Hilbert curve serialization in state-of-the-art point cloud architectures.


% This paper provides a theoretical framework to analyze the generalization error of these approaches by bounding their covering numbers. We establish a strict generalization hierarchy: equivariant models and group averaging achieve error bounds that are better than or equal to those of models using canonization, which in turn are better than or equal to models that ignore invariance completely. Crucially, we introduce a method to compare competing canonizations, defining them as "optimal" or "poor" based on their proximity to the covering number of the quotient space. We give several examples of optimal and poor canonizations. Finally, applying this to permutation groups for point cloud processing, we prove that the covering number of lexicographical sorting is far higher than the covering number attained by Hilbert curve canonization. Experimentally, we provide empirical evidence for our conclusions, showing the advantages of group averaging over canonization and of Hilbert canonization over lexicographic sorting. 
\end{abstract}
\section{Introduction}
The integration of geometric priors into neural networks has become a cornerstone of modern machine learning, particularly for applications involving graphs, 3D point clouds, and molecular structures. In these domains, the underlying data often exhibits inherent symmetries, meaning the target function to be learned is equivariant to a specific group action. Exploiting these symmetries is known to improve sample complexity and generalization \cite{elesedy2021provably,brehmer2025doesequivariancematterscale}.

While specialized architectures, which are equivariant by design, are a common method for processing symmetric data, there is growing interest in achieving equivariance by applying generalized group averaging methods to non-equivariant backbones. This family of methods includes full group averaging, which is theoretically sound but intractable for all but very small groups, group augmentation, which can be seen as an efficient approximation of group averaging, frame averaging \cite{punyframe}, which enables equivariant averaging over subsets of the group, and canonization, which maps every group orbit to a single, consistent "canonical" element. Canonization is the most efficient of all the generalized group averaging, as it only involves processing a single element. At the same time, it enjoys the same universal approximation guarantees as full group averaging \cite{kaba2023equivariance}. These observations have motivated researchers to define canonizations in several diverse domains. These include sign canonizations for spectral embeddings of graphs \cite{ma2023laplacian,ma2024a,hordan2025spectral}, canonizations of graphs \cite{lin2024equivariance} and point clouds \cite{kaba2023equivariance,pmlr-v235-baker24a,friedmann2025,zhou2026rethinkingdiffusionmodelssymmetries} including the celebrated Point Transformer V3 \cite{wu2024ptv3}, and canonization of the SMILES representation of small molecules \cite{canonicalSMILES}.



In this paper, we provide a theoretical study of canonization from the perspective of generalization. Our focus is both on comparing canonization with other options for enforcing equivariance and on comparing the quality of different canonizations. Our main  contributions are threefold:

\begin{enumerate}
\item \textbf{Canonization vs. Group Averaging:} We prove that the generalization bounds of canonized models are at best equal to the bounds of group-averaged models, and at worst equal to the bounds of non-equivariant models.

\item \textbf{Canonizations and continuity:} We show that continuous, isometric canonizations attain the \emph{optimal} generalization error achieved by group averaging, while discontinuous canonizations can lead to \emph{poor} canonizations whose error bounds are identical to non-equivariant models.

\item \textbf{Point Clouds Serialization:} We prove that 
canonizing point clouds with respect to permutation via lexicographical sorting is inferior to the  Hilbert canonizations used in  Point Transformer V3 \cite{wu2024ptv3}, thus establishing a theoretical justification for the empirical success of this method.
\end{enumerate}


\subsection{Related Work} 
\paragraph{Canonization vs. Group Averaging} It was observed that randomized SMILES can often outperform canonized SMILES \cite{ArusPous2019,Bjerrum2017SMILESEA} (see also \cite{ito2026random}). Our first conclusion gives a theoretical justification for this empirical observation.

 \paragraph{Discontinuity} It was shown in \cite{dym2024equivariant,pmlr-v235-baker24a} that in many equivariant learning scenarios, continuous canonizations are mathematically impossible. Attempts to mitigate this issue were made in  \cite{tahmasebi2025regularity,lin2026adaptive}. Our second conclusion complements these works by showing how discontinuity of canonization hurts the generalization of the canonized model.  

\paragraph{Equivariance and Generalization} Many papers have considered generalization analysis of equivariant models \cite{vasileiou2025survey,petrache2023approximationgeneralization,maskey2025generalization, franks2024weisfeiler}, but they do not consider canonizations. To our knowledge, the only paper considering the generalization of canonizations is \cite{tahmasebi2025generalization}. They prove that in some setting the \emph{approximation error}  of canonized models is lower than group averaged models, and as a result the \emph{expected error} of canonized models can be lower when the model is presented with enough samples. In contrast, our first conclusion is that the \emph{generalization error} of group average models is always lower than the generalization error of canonized models.

% \paragraph{Relation to \cite{tahmasebi2025generalization}}
% We note that this conclusion differs from the analysis of \cite{tahmasebi2025generalization}, who concluded that in some settings, canonization can generalize better than group averaging, which is an example of an invariant model. The reason for this difference is that their analysis focuses on reducing the expected loss, rather than the generalization error as in our analysis. Indeed, our analysis does not rule out the possibility that a canonized model will be more expressive than the invariant model, and as a result, the expected loss of the canonized model will be lower than the invariant model, despite the generalization error being larger. However, we believe that such expressive differences are limited to the relatively simple function classes considered in their paper, as when universal models with a large number of parameters are considered, the expressive power of group averaging and canonization is nearly identical, rendering this effect negligible. 

\subsection{Notation}
Throughout the paper, we will consider a metric space $(K,\rho) $, where $K$ is endowed by the action of a group $\G$. We will assume that the metric is $\G$ invariant, i.e., $ \rho(g\cdot x, g \cdot y)=\rho(x,y)$ for all $x,y\in K$ and $g \in \G $.  

\paragraph{Quotient space} We denote the orbit of $x\in K$ by  $[x]=\{gx| \quad g\in G \}$, and denote the quotient space to be the space of orbits $K/\G:=\{[x]|x \in K\}$. On this space, we define a metric $\rho_\G([x],[y]):=min_{g \in \G}\rho(g\cdot x,y)$. The $\G$ invariance of $\rho$ guarantees that $\rho_{\G}$ is symmetric and satisfies the triangle inequality. To ensure that $\rho_{\G}([x],[y])=0 $ implies that $[x]=[y]$, we explicitly add the assumption that the minimum in the definition of $\rho_\G$ always exists. A triplet $(K,\rho,\G)$ satisfying the conditions of the last two paragraphs is called a \emph{module}, and all examples considered in this paper are modules.

\textbf{Covering number}
Let $(K,\rho)$ be a metric space. We say that $C \subset K$ is an $\epsilon$ cover of $K$,  if $\forall x \in K, \exists y \in C: \rho(x,y)\leq  \epsilon$. The minimal $N$ such that there exists an $\epsilon$ cover of cardinality $N$ is called the $\epsilon$ covering number of $K$, and is denoted by $\N(K,\rho,\epsilon)$. 

\paragraph{Canonization} 
We say that $c:K \rightarrow K$ is a canonization if for every $x\in K$, (i) $x$ is mapped to an element in its orbit, namely $c(x)\in [x]$, and (ii) all elements in the orbit $[x]$ are mapped to the same element. 

\section{Generalization, Canonization and Group Averaging}\label{sec:gen}
Consider a supervised learning problem of learning an unknown function $f:X\to Y $ from a finite set of $n$ samples $S=\{(x_i,f(x_i)| \quad i=1,\ldots,n) \}$ of the function. Assume that we have some algorithm (e.g., gradient descent) which, given $S$, returns a hypothesis $h_S:X \to Y $. Let $\ell:Y \times Y \to \RR  $ be a function which will be used as a loss function. The \emph{expected loss} $\mathcal{L}$ and the \emph{empirical loss} $\Lemp $ are defined via 
$$\mathcal{L}_{exp}(h_S) \triangleq \mathbb{E}_{x \sim \mu} \ell(h_S(x),f(x)), \quad \Lemp(h_S)\triangleq\frac{1}{n}\sum_{i=1}^n \ell(h_S(x_i),f(x_i))$$

The \emph{generalization} error of $h_S$ is the absolute difference between the expected and empirical loss. The generalization error of $h_S$ can be bounded using the Lipschitz constants of the function  $h_S,f,\ell$ and the covering number of the domain $X$, as follows:

\begin{restatable}{theorem}{MainGeneralizationBound}
\label{thm:MainGeneralizationBound}
Let $(X,\rho_X) $ and $(Y,\rho_Y) $ be  metric spaces and assume $X$ is compact. Let $f:X\to Y $ be a $c_f$  Lipschitz function, and let $\ell:Y\times Y \to [0,M] $ be a $c_\ell $ Lipschitz function. Then for any $\epsilon,\delta>0$ and natural $n$,  for a training sample set $S$ generated by $n$ IID draws from a distribution $\mu$,  with probability of at least $ 1-\delta$, we have 
$$|\mathcal{L}_{exp}(h_S) - \Lemp(h_S)| \le 2c_\ell \max\{1,c_{h_S} \} (1+c_f)\epsilon+M \sqrt{\frac{2\textcolor{blue}{\mathcal{N}( X, \rho_X,\epsilon)}\ln 2 + 2 \ln(1/\delta)}{n}}$$
where $c_{h_S}$ denotes the Lipschitz constant of $h_S$.
\end{restatable}

This theorem follows easily from a well-known theorem from \cite{xu_robustness_2012}, which is often employed in the study of neural network generalization  \cite{vasileiou2025survey}.  The derivation from the Theorem is given in Appendix \ref{app:gen}.  

We now apply this theorem in the setting where the learned function $f$ is invariant to a group action. In the following proposition, we consider three types of models: (i) models that do not exploit function symmetry at all, (ii) models that are invariant by construction, and (iii) non-invariant models composed with canonizations. 

\begin{restatable}{proposition}{generalizationbound}
\label{propGen}
Let $(K,\rho,\G)$ be a module. Let  $(Y,\rho_Y) $ be a metric space, and assume $K$ is compact. Let $f:K\to Y $ be a $c_f$  Lipschitz function \textbf{which is $\G $ invariant}, and let $\ell:Y\times Y \to [0,M] $ be a $c_\ell $ Lipschitz function , then for any $\epsilon,\delta>0$ and natural $n$,  for a training sample set $S$ generated by $n$ IID draws from the distribution $\mu$,  with probability of at least $ 1-\delta$, we have 
$$|\mathcal{L}_{exp}(h_S) - \Lemp(h_S)| \le 2c_\ell \max\{1,c_{h_S} \} (1+c_f)\epsilon+M \sqrt{\frac{2\textcolor{blue}{\N(h_S)} \ln 2 + 2 \ln(1/\delta)}{n}}$$

where 
\begin{equation}\label{eq:differentN}
\textcolor{blue}{\N(h_S)}=\threepartdef{\N(K,\rho,\epsilon)}{h_s \text{ is } c_{h_s} \text{ Lipschitz }}{\N(c(K),\rho,\epsilon)}{h_s=\tilde{h}_s\circ c, \tilde{h}_s \text { is } c_{h_s} \text{ Lipschitz, and }  c \text{ is a canon.}} {\N(K/\G,\rho_{\G},\epsilon)}{h_s \text{ is } c_{h_s} \text{ Lipschitz and  } \G \text{ invariant }} 
\end{equation}
\end{restatable}

\begin{proof}[Proof idea]
The claim for a Lipschitz function $h_s$ without any invariant structure follows immediately from Theorem \ref{thm:MainGeneralizationBound}. For a function of the form $h_s=\tilde h_s \circ c $, the claim follows from Theorem \ref{thm:MainGeneralizationBound} when replacing the domain $K$ with $c(K)$ and the function $h_s$ with $\tilde h_s$. The claim for $h_s$ which is both Lipschitz and invariant follows from the fact that due to invariance $h_s, f$ can be identified with function $\hat h_s, \hat f : K/\G \to Y $ satisfying $\hat h_s([x])=h_s(x), \hat f([x])=f(x), \forall x\in K $, and the Lipschitz constant remains unchanged after the identification (see Lemma 20 in \cite{siegel2026quantitativeapproximationratesgroup}).
\end{proof}
\nd{to self: add proof for equivariant setting in appendix, mention in main text}

For each of the three model types considered in Proposition \ref{propGen}, a different value of $\textcolor{blue}{\N(h_S)}$ is provided. The following simple lemma establishes the relationship between these three values: 

\begin{lemma}
    For any module $(K,\rho,\G)$, and any $\epsilon>0$,  the following inequality holds: 
    \begin{equation}\label{eq:coveringIneq}
        \N(K/\G,\rho_{\G},\epsilon)\leq \N(c(K),\rho,\epsilon)\leq \N(K,\rho,\epsilon)
    \end{equation}
\end{lemma}




\begin{proof}
The inequality $\N(c(K),\rho,\epsilon)\leq \N(K,\rho,\epsilon)$ is immediate since $c(K)\subseteq K$. For the other inequality, note that the quotient mapping $q:c(K)\to K/\G $ defined by $q(x)=[x]$ is onto and $1$-Lipschitz, and therefore it transforms any $\epsilon$ cover of $c(K)$ to an $\epsilon$ cover of $K/\G$, thus implying that $\N(K/\G,\rho_{\G},\epsilon)\leq \N(c(K),\rho,\epsilon) $.

% For the other inequality, let $C$ be an $\epsilon$-cover of $c(K)$ with $|C|=\N(c(K),d,\epsilon)$, and define
% \[
% \widehat{C}:=\{[x]\mid x\in C\}.
% \]
% Clearly $|\widehat{C}| \leq|C|$.

% We claim that $\widehat{C}$ is an $\epsilon$-cover of $K/\G$ under $\rho{\G}$.
% Let $[y]\in K/\G$. Since $c(y)\in c(K)$ and $C$ covers $c(K)$, there exists $x\in C$ such that
% $\rho(x,c(y))\le \epsilon$.
% Using that $[c(y)]=[y]$, we obtain
% \[
% \rho{\G}([x],[y]) = \rho{\G}([x],[c(y)]) \le \rho(x,c(y)) \le \epsilon.
% \]
% Thus $\widehat{C}$ covers $K/\G$, and therefore
% $\N(K/\G,\rho{\G},\epsilon)\le |\widehat{C}|\le |C|=\N(c(K),d,\epsilon)$.
\end{proof}
Informally, the conclusion we draw from this lemma is as follows:
\begin{conclusionbox}{Conclusion 1}
 The generalization bounds of canonized models are at best equal to the bounds of group-averaged models, and at worst equal to the bounds of non-equivariant models.
\end{conclusionbox}

\bgf{isnt the proof done for invariant models?
In general, it feels "equivariant" and "invariant" are switched in and out, which is a bit confusing. Maybe clarify? }

While this conclusion is valid for all equivariant models,  it is most relevant for equivariant models obtained from group averaging a non-equivariant model. In  this setting, the same backbone model can be used for canonization and group averaging, and so a direct comparison can be made. 
 
We conclude this section by showing that the ratio between the largest and smallest covering numbers in \eqref{eq:coveringIneq} is bounded by the group's cardinality:
\begin{restatable}{proposition}{cardinality}[See proof in ~\ref{upper_bound}]
\label{prop:cardinality}
Let $(K,\rho,\G)$ be a module, and assume $\G$ is finite. Then
\[
\N(K,\rho,\epsilon)\le |\G|\cdot \N(K/\G,\rho_{\G},\epsilon).
\]
\end{restatable}


\begin{figure}[htbp]
    \centering
     % \caption{\textbf{Comparison of canonizations for the action of $\G=\{-1,1\}$ on $\mathbb{R}$.} \textit{Left:} The isometric canonization $c(x)=|x|$. Because it is continuous, it is an optimal canonization that achieves the quotient space covering number. \textit{Middle:} The canonization $c_1$, which exhibits isolated discontinuities at $x = \pm 0.5$ (indicated by the open and closed limits). \textit{Right:} The extreme canonization $c_\infty$, which shatters the domain across both branches based on rationality. As established in Proposition \ref{prop:poor}, the partial limits in the discontinuous canonizations $c_1$ and $c_\infty$ allow for domains where the covering numbers hit the non-equivariant upper bound, rendering them poor canonizations.}
      \caption{\it Visualization of three canonizations for the action of $\G=\{-1,1\}$ on $\mathbb{R}$. The natural  canonization $c(x)=|x|$, which is shown to be optimal (Example \ref{prop:isometries}), and two discontinuous canonizations, which on an appropriate domain are poor (Example \ref{ex:sign_poor}).}
    \label{fig:1d_canonizations}
    \includegraphics[width=0.8\linewidth]{canonization_neurips_2027.pdf}
   
\end{figure}


\section{ Canonizations and Continuity}\label{sec:evaluation}
In the previous section, we established upper and lower bounds for the covering numbers of canonizations. Our focus in this section is to study when these bounds are achieved. As we will see, this will be related to the continuity of the canonization.  We begin by defining optimal and poor canonizations. 
\begin{definition}
Let $(K,\rho,\G)$ be a module and let $\epsilon$ be a positive number. We say that $c$ is $(K,\epsilon) $ \emph{optimal} if it attains the lower bound in \eqref{eq:coveringIneq}, namely 
$
\N(c(K),\rho,\epsilon) = \N(K/\G,\rho_{\G},\epsilon).
$
We say that $c$ is $(K,\epsilon) $ \emph{poor}  if it attains the upper bound in \eqref{eq:coveringIneq}, namely 
$
\N(c(K),\rho,\epsilon)=\N(K,\rho,\epsilon).
$
\end{definition}
We will next show conditions for optimal and poor canonizations, and show how these conditions relate to the continuity of the canonization. 

\textbf{Optimal canonizations:}
First, we prove canonizations, which are isometries, are optimal. We say that a canonization $c:K\to K $ is an \emph{isometry} if $\rho(c(x),c(y))=\rho_{\G}([x],[y])  $ for all $x,y\in K$.
\begin{restatable}{claim}{isometry}\label{claim:isometry}[proof in \ref{Proof_isometry}]
    Any isometric canonization $c:K\to K$ is $(K,\epsilon) $ optimal for all $\epsilon>0$.
\end{restatable}

We next use this claim to give several examples of  optimal canonizations:
\begin{restatable}{example}{isometries}\label{prop:isometries}
    The following canonizations are isometries (see proof in \ref{Proof_isometries})
\begin{enumerate}
\item The canon. $c(x)=|x|$ with respect to the action of $\{-1,1\} $ on $\RR$ and the standard metric.
\item The canon. $\mathrm{sort}:\RR^n \to \RR^n $ with respect to the action of permutations and a metric $\rho$ induced by a $p$ norm.
\item The centralization canon. defined for point clouds in $\RR^{d\times n}$ with respect to the action of translation by vectors in $\RR^d$, with the metric $\rho$ induced by the Frobenius norm.
\end{enumerate}
\end{restatable}


%  In \cite{dym2024equivariant}, it is shown that in many equivariant learning settings, a continuous canonization (and thus in particular an isometric canonization) is not available. However, the possible implications of this for learning are not explained.
\textbf{Poor canonizations:} We next show that discontinuities in canonizations can lead to poor canonizations. This occurs under two conditions: firstly, the canonization is 'strongly discontinuous' in the sense that it has $|G|$ partial limits. Secondly, the domain $K$ of the group action is chosen adversarially to target the canonization's discontinuity. 
\begin{restatable}{proposition}{poorcanonization}[proof in \ref{poor_proof}]
\label{prop:poor}
    Let $(V,\rho,\G)$ be a module, where $\G$ is a finite group. Let $c$ be a canonization such that $c$ has $|G|$ partial limits at a point $x\in V$, then for any small enough $\epsilon>0$ there exists a $\G$ invariant set $K\subseteq V$ such that $\N(c(K),\rho,\epsilon)=\N(K,\rho,\epsilon)=|\G| \cdot \N(K/\G,\rho,\epsilon)$. 
\end{restatable}


Based on this proposition, we give two examples of poor canonizations:

\begin{example}\label{ex:sign_poor}
Consider the action of $\G=\{-1,1\}$ on $V=\RR $. Consider the canonizations
$$c_1(t)=\twopartdef{|t|}{|t|>1/2}{-|t|}{|t|\leq 1/2}, \quad  c_\infty(t)=\twopartdef{|t|}{t\in \mathbb{Q}}{-|t|}{t\not \in \mathbb{Q}} .$$
Both of these canonizations are discontinuous and have a point with $|G|=2 $ partial limits. Therefore, by  Proposition \ref{prop:poor} for every small enough $\epsilon>0$ we can choose an adversarial domain for which the canonization is poor. For $c_1$ we can choose for every $\epsilon\in (0,1)$ the set  $K=B_\epsilon(-\frac{1}{2})\cup B_\epsilon(\frac{1}{2})$ and then $\N(K,\rho,\epsilon)=\N(c_1(K),\rho,\epsilon)=2$ and so the canonization is poor. %This example stresses the fact that the choice of  $K$ depends on $\epsilon$. In fact, for any given $K$, it is not difficult to see that the covering number of $c_1(K)$ is larger than the covering number of $K/\G $ by at most one, and therefore we obtain that the canonization is optimal at the limit $\epsilon \rightarrow 0 $, in the sense that $\lim_{\epsilon \rightarrow 0}\frac{\N(c(K),\rho,\epsilon)}{\N(K/\
%G,\rho_\G,\epsilon)}=1$.
In the canonization $c_\infty$ the discontinuities are much more extreme. In this case, we can choose a non-adversarial domain $K=[-1,1]$, and the canonization will be poor for any $\epsilon>0$ because the covering number of a set and its closure are the same, and in our case, the closure of $c_{\infty}(K)$ will be all of $K$. 
\end{example}

%In the last example, the group cardinality was 2, and as a result, the difference in covering numbers between poor and optimal canonizations is not dramatic. In the next example, we will consider large groups where the differences are much more substantial. 

\begin{example}\label{ex:multisort}
Consider the action of the permutation group $\G=S_n$ on the space of matrices $\RR^{d\times n} $ by permuting the columns of the matrices. We call such matrices 'point clouds'. We assume that $n,d\geq 2$ (the case $d=1$ sorting is an isometry by Example \ref{prop:isometries}). A natural canonization in this scenario is $\csort$, which permutes the columns of the matrices so that the first row of the matrix is sorted from small to large, and in the event of ties in the first row, sorts according to the second row, and continues in this fashion until a permutation is uniquely defined. The canonization $\csort$ is discontinuous. Moreover, the canonization has $n!=|\G|$ partial limits at any matrix  $M\in \RR^{d \times n} $ satisfying that the matrix does not contain two identical columns, but $M_{1,1}=M_{1,2}=\ldots=M_{1,n} $. Thus, according to Proposition \ref{prop:poor}, for every small enough $\epsilon>0$ there exists an adversarial domain for which the canonization is poor. 
\end{example}
\begin{conclusionbox}{Conclusion 2}
We show that isometric canonizations are optimal, while discontinuous canonizations can yield poor canonizations. These results give a rigorous justification for the importance of continuous canonizations.
\end{conclusionbox}


\section{Case study: Permutation Canonizations}\label{sec:case_study}

In Example \ref{ex:multisort}, we showed that lexicographical sorting is a poor canonization on an adversarial domain. In this section, we analyze the covering numbers of lexicographical sorting when considering a fixed, 'natural' domain of point clouds. We will show that even in this non-adversarial setting, the covering number of lexsort canonization is much higher than the covering number of the quotient space: we show this by showing that the covering number of lexsort grows exponentially in $n$, the number of points in the point cloud, while the covering number of the quotient space only grows polynomially in $n$.   Moreover, we will consider an alternative canonization based on Hilbert curves, and show that its covering numbers also grow only polynomially in $n$. The significance of this result is that it gives a theoretical justification for the choice of  Point Transformer V3 \cite{wu2024ptv3} to use Hilbert canonizations for point cloud serialization.  

Throughout this section we consider the module $(K,\rho,\G) $ with the domain $K=[0,1]^{d\times n}$, the distance induced by the element-wise infinity norm $\rho(X,Y)=\|X-Y\|_\infty $, and the permutation group $\G=S_n$. We  will show that for fixed $\epsilon,d$, the covering number of $\csort$ grows exponentially in $n$, while the covering number of the quotient space and Hilbert canonizations only grows polynomially in $n$. We begin by proving the exponential growth of lexsort canonizations:

\begin{restatable}{lemma}{lexsortcovering}[proof in \ref{lexsortcovering}]
\label{lem:lexsortcovering}
Let $k,d,n\geq 2$ be natural numbers, and set $\epsilon=1/(2k), K=[0,1]^{d\times n}$ and $\G=S_n$.    Then the covering number of the lexicographical sorting canonization satisfies 
\begin{equation}\label{eq:exp}\left( \frac{1}{2\epsilon}\right)^{(d-1)\cdot n + 1}\leq \N(\csort(K),\rho,\epsilon).\end{equation}
\end{restatable}
\begin{proof}[Proof idea]
The lexsort canonization is ineffective on the set of its 'severe discontinuities'- point clouds whose first row is constant. These point clouds form a $(d-1)\cdot n+1 $ dimensional hypercube. The covering number of $\csort(K)$ is upper bounded by the covering number of this hypercube, which is the left-hand side of \eqref{eq:exp}.
\end{proof}
Next, we prove the polynomial growth of the quotient space's covering number:
\begin{restatable}{lemma}{quotientcovering}
\label{lem:quotient_covering}
Let $n,d$ be natural numbers and $\epsilon\in (0,1)$. Set $K=[0,1]^{d\times n}$ and $\G=S_n$. Then the covering number of the quotient space satisfies
\begin{equation}\label{eq:combinatorial}
\N(K/\G,\rho_{\G},\epsilon)
\leq
\binom{n+k^{d}-1}{n},
\text{ where }
k=\left\lceil \frac{1}{2\epsilon} \right\rceil .
\end{equation}
\end{restatable}
\begin{proof}
Set $k=\lceil 1/(2\epsilon) \rceil $. The set $K$ can be covered by $k^d$ hypercubes of diameter $k$ and radius $1/(2k)\leq \epsilon $. Let $C$ denote the collection of centers of these hypercubes, and after applying the quotient map  $\pi(c) = [c]$ we obtain points $\pi(C)$ in the quotient space which are an $\epsilon$ cover of $K/\G$. Note that points in $C$ which are permutations of each other are mapped by $\pi$ to the same orbit. Therefore, the cardinality of the cover $\pi(C)$ is the number of equivalence classes with respect to the relation on $C$ defined by the group:  $c_1 \equiv c_2 \iff \exists \pi \in S_n : c_1 = \pi(c_2)$. As explained in more detail in the proof, the number of equivalence classes is the combinatorial expression in \eqref{eq:combinatorial}.  %\begin{equation*}
%|\pi(C)|= \binom{n+k^d-1}{n}
%\end{equation*} 
\end{proof}



% \begin{lemma}
%     Assume $d = l_{\infty}$, so $\N(K,d,\epsilon)= (\frac{1}{2\cdot \epsilon})^{nd}$
% \end{lemma}
% \begin{proof}
%     It's easy to see that each line of length $1$, we can cover using $\frac{1}{2\cdot \epsilon}$ $\epsilon$ size lines, so using the product to cover $K$, we need $\frac{1}{(2\epsilon)^{nd}}$ balls to cover it.
%     On the other hand denote by $F$ some other covering. 
%     Note that $K \subseteq \cup_{c \in F} B(c, \epsilon)$
%    So $1=\mu(K)\leq \mu(\cup_{c \in F} B(c, \epsilon))\leq \sum_{c \in F} \mu(B(c, \epsilon))=|F|\cdot (2\epsilon)^{n \cdot d}$ so $\frac{1}{(2\cdot \epsilon)^{nd}}\leq |F|$. 
%    Thus, we got the lower bound and the upper bound as desired.
% \end{proof}


% \begin{conclude}
% We obtain an upper bound for the covering number of $K/\G $ and a lower bound for the covering number of $\csort(K) $. The upper bound grows polynomial in $n$ while the lower bound grows exponentially in $n$, showing the large gap between covering numbers as $n\rightarrow \infty$ and that in this setting the canonization is far from being optimal. We next introduce Hilbert canonizations and show that they do exhibit polynomial growth. 
% \end{conclude} 
\subsection{Hilbert Curve Canonization}
Finally, we will give an intuitive  explanation of the Hilbert canonization and formulate the claim that its covering number exhibits only polynomial growth in the point cloud dimensionality. For a full definition of the canonization and proof see Appendix \ref{app:hilbert}.
\begin{figure}[h]
    \centering
        \caption{\it The Hilbert curves $H_m$ for $d=2$ and $m=1,2,3$ induce an ordered on a grid in $[0,1]^d $ with $2^{md}$ points. On the right we see the discontinuous lexsort ordering on the grid.}
    \includegraphics[width=1.1\linewidth]{figs/hilbert_with_lex_sort.pdf}
    \label{fig:hilbert}
\end{figure}

The Hilbert curve is a space-filling curve, mapping the unit interval $[0,1]$ continuously onto a hypercube $[0,1]^d $. It is defined as a limit of curves $H_m:[0,1]\to [0,1]^d $ which are not onto but whose image becomes more and more 'dense' in $[0,1]^d $ as $m$ grows, as shown in Figure \ref{fig:hilbert} for the case where $d=2$. The curves $H_m$ can be used to define an order on $[0,1]^d$ (at least on points in the image of the curve) and thus induce a canonization.  Intuitively, the advantage of this order over the lexsort ordering is that this ordering is continuous (see Figure \ref{fig:hilbert}) again. However, while previous work \cite{wu2024ptv3,hilbertnet} cited this continuity as justification for  Hilbert ordering, a formal justification of why this continuity is important for learning has not yet been provided. We address this caveat by showing the polynomial growth of the covering number of Hilbert canonizations:
\begin{restatable}{theorem}{hilbert}\label{thm:hilbert}
Let $d,n,m\geq 2$ be natural numbers, let $\epsilon\in (0,1)$, and assume that $\epsilon>2^{-m-1} $ set $ K=[0,1]^{d\times n}$ and $\G=S_n$.    Then the covering number of the Hilbert canonization $c_m$ satisfies 
$$\N(c_m(K),\rho_\infty,\epsilon)\leq  \binom{n+\lceil 1/ (2\delta) \rceil-1}{n},   \text{ where } \delta= \frac{\left(\epsilon-2^{-m-1}\right)^d}{4}$$
\end{restatable}
\begin{proof}[Proof idea]
When restricting to $n$-tuples in $K$ which are all in the image of $H_m$, the Hilbert canonization can be written as $H_m \circ \mathrm{sort} \circ H_m^{-1} $. Since $\mathrm{sort}$ (Proposition \ref{prop:isometries}) is an isometry, the image of $  \mathrm{sort} \circ H_m^{-1}$  can be covered efficiently. This cover can then be pushed forward to cover the image of   $H_m\circ \mathrm{sort} \circ H_m^{-1}$, using the fact that $H_m$ is Holder continuous.
% The proof is based on three observations: (i) an $\epsilon$ cover of $[0,1]^{d\times n} $ can be obtained from an $\epsilon-2^{-m-1}$ cover of $\left(H_m[0,1]\right)^n $. (ii) $H_m$ is Holder continuous with known constants, for an appropriate $\epsilon'$, an $\epsilon'$ cover of $[0,1]^n $ can be pushed forward by $H_m$ to an $\epsilon$ cover of $\left(H_m[0,1]\right)^n $. (iii) 
\end{proof}

 Theorem \ref{thm:hilbert} shows that the covering number of the Hilbert canonization grows at worst polynomially in $n$, for fixed $\epsilon>0$. This is in contrast with the exponential growth of the sorting canonization, proving a substantial gap between them. 

\begin{conclusionbox}{Conclusion 3}
 Our analysis proves that the covering number of Hilbert canonizations is far superior to lexsort canonization, giving a theoretical justification for its superior performance in practice. 
\end{conclusionbox}


To obtain more intuition for the gap between our bounds for the covering numbers, in Table \ref{tab:bounds} we present bounds for the covering number of the quotient space, the Hilbert canonization and the sorting canonization, as well as the covering number of the hypercube $[0,1]^{d\times n}$ which is known to be $k^{nd}$ when $\epsilon=1/(2k)$ (see proof in \ref{covering_cube}). We take characteristic values encountered in point cloud learning scenarios:  point dimensions are $d=3$, a very modest value of $\epsilon=1/6$, and $n$ varying between $250$ and $2000$. 

\renewcommand{\arraystretch}{1.3} % More vertical space for large exponents
\begin{table}[h]
\centering
\label{tab:bounds}
\caption{\it Comparison of covering number bounds as a function $n$ between the quotient space,  the Hilbert and Lexsort canonizations, and the hypercube}
\begin{tabular}{lccccc}
\hline
 \textbf{Covering number vs. $n$} & $\mathbf{n=250}$ & $\mathbf{n=500}$ & $\mathbf{n=750}$ & $\mathbf{n=1000}$& $\mathbf{n=2000}$ \\ \hline
Quotient (upper bound) & $2.1 \times 10^{36}$  & $7.4 \times 10^{43}$  & $2.2 \times 10^{48}$  & $3.5 \times 10^{51}$ & $2.0 \times 10^{59}$  \\
Hilbert (upper bound) & $5.3 \times 10^{193}$ & $7.9 \times 10^{278}$ & $5.0 \times 10^{336}$ & $5.0 \times 10^{380}$& $4.4\times 10^{494}$ \\
Lexsort (lower bound) & $1.1 \times 10^{239}$ & $4.0 \times 10^{477}$ & $1.4 \times 10^{716}$ & $5.2 \times 10^{954}$ & $9.2\times 10^{1908}$ \\
Hypercube (exact) & $6.9\times 10^{357}$ & $4.8\times 10^{715}$ & $3.3\times 10^{1073}$ & $2.3\times 10^{1431}$ & $5.3 \times 10^{2862}$\\
 \hline
\end{tabular}
\end{table}
We note that our bounds for the quotient space and Hilbert canonizations are upper bounds, while the Lexsort bound is a lower bound. Thus, the gap between the covering number of Lexsort and that of the other methods may be even larger than shown in the table. We also note that despite the conservative choice of $\epsilon=1/6$, all covering numbers in the table are enormous. This picture may be over-pessimistic, as the 'true data distribution' will be supported on a 'small' set $K\subseteq [0,1]^{d\times n} $ of 'realistic point clouds'. An experiment estimating the covering number on the Modelnet40 dataset will be shown later on in Table \ref{tab:distances_both}.











\section{Experiments}
We conduct several experiments to verify our theoretical insights and their manifestation in practical learning scenarios. 

\paragraph{Covering Number on ModelNet} In our first experiment, we considered the ModelNet \cite{wu20153d} point cloud classification task. Our first goal was to see whether the covering number hierarchy between the four different methods considered in Table \ref{tab:bounds} is preserved when replacing the large domain $[0,1]^{3\times n}$ with the 'ModelNet point-cloud domain' $K$. Namely, our interest is in estimating the covering number of $K $, its image under the lexsort and Hilbert canonizations, and the covering number of the quotient space $K/\G $.    

We do not have access to the 'true point cloud' domain $K$, but only to the train and test sets. Thus, to estimate the covering number, we compute, for each test sample $s$, the minimum distance from $s$ to all training samples (for computational efficiency, we consider only samples with the same label as $s$). We call this number $q_s$. We can then estimate the covering number by taking $\epsilon=\max_{s\in \mathrm{Test}} q_s $, as this implies that the train set is an $\epsilon$-cover of the test set. We call this measure the \emph{maximal coverage}. We also consider a more robust measure, the \emph{average coverage} $\mathrm{mean} \{ q_s| s\in \mathrm{Test} \} $. The results, shown in Table \ref{tab:distances_both}, show that the expected covering number hierarchy is also preserved when considering real data. 

\begin{table}[h]
\centering
\caption{\it Distance from each test sample to its nearest train sample for $n=256$ on ModelNet10 and ModelNet40, reporting both mean and max over the test set. Lower is better.}
\label{tab:distances_both}
\begin{tabular}{lcccc}
\toprule
 & \multicolumn{2}{c}{\textbf{Mean coverage} $\downarrow$} 
 & \multicolumn{2}{c}{\textbf{Max coverage} $\downarrow$} \\
\textbf{Method} 
& \textbf{ModelNet10} 
& \textbf{ModelNet40} 
& \textbf{ModelNet10} 
& \textbf{ModelNet40} \\
\midrule
Regular distance & 0.4545 & 0.3821 & 0.629  & 0.7304 \\
Sort distance    & 0.3770 & 0.2988 & 0.5407 & 0.5550 \\
Hilbert distance & 0.2013 & 0.1955 & 0.5283 & 0.5490 \\
Group distance   & \textbf{0.0917} & \textbf{0.0957} & \textbf{0.2742} & \textbf{0.4880} \\
\bottomrule
\end{tabular}
\end{table}

\paragraph{Classification on ModelNet} Next, we consider whether the hierarchy in covering numbers between the different methods translates into a corresponding hierarchy in accuracy and generalization error. We choose a simple, purely non-invariant architecture: a Global MLP, and compare the performance of this network (i) without canonization, (ii) with lexsort canonization, and (iii) with Hilbert canonization (we cannot implement group averaging here since the group is very large). The results, shown in Table \ref{tab:best_modelnet10and40}, reveal that the expected hierarchy is obtained, with Hilbert performing best both on ModelNet40 and ModelNet10, followed by the lexsort canonization and a plain MLP. The Generalization error column (defined as Test acc.-Train acc.) indicates that this performance gap is due to improved generalization.  

\begin{table}[h]
\centering
\caption{\it On the ModelNet classification task, we compare vanilla MLP architectures in three configurations: (a) no canon. (b) lexsort canon. and (c) Hilbert canon. As predicted by our theory, the Hilbert canon. leads to better generalization, which translates into improved overall accuracy.}
\label{tab:best_modelnet10and40}
\begin{tabular}{lcccc}
\hline
 & \multicolumn{2}{c}{ModelNet10} & \multicolumn{2}{c}{ModelNet40} \\
Model & Test Acc $\uparrow$ & Gen. error $\downarrow$ & Test Acc $\uparrow$ & Gen. error $\downarrow$ \\
\hline
MLP (no canon.)  & 0.57 & 0.43 & 0.442 & 0.558 \\
MLP (lexsort)    & 0.755 & 0.226 & 0.693 & 0.2556 \\
MLP (Hilbert)    & \textbf{0.7815} & \textbf{0.1167} & \textbf{0.74} & \textbf{0.0861} \\
\hline
\end{tabular}
\end{table}

To further verify the generality of our findings, we also trained the three alternatives on smaller samples from the ModelNet40 dataset. We find that the $\text{Hilbert}>\text{lexsort}>\text{no canon.}$ hierarchy is preserved in all experiments, as shown in Table \ref{tab:data_scarcity}. 
We note that this experiment is only illustrative; the results attained by simple vanilla MLPs, even with Hilbert canonization, are far from state-of-the-art. In the context of our theory, this can also be explained through the fact that successful models for this task are typically permutation invariant, and hence their generalization is governed by the covering number of the quotient space, which is always better than the covering number of canonizations (Conclusion 1).


% \begin{table}[h]
% \centering
% \caption{On the Modelnet pointcloud classification task, we comparing the performance of vanilla MLP architectures in three configurations: (a) no canon. (b) sorting canon. and (c) Hilbert canon. As predicted by our theory Hilber canon. leads to better generalization, which translates into improved overall accuracy.}
% \label{tab:best_modelnet10and40}
% \begin{tabular}{lcccccc}
% \hline
%  & \multicolumn{3}{c}{Modelnet10} & \multicolumn{3}{c}{Modelnet40} \\
% %\cline{1-7}%\cline{2-4} \cline{5-7} 
% Model & Train Acc & Test Acc & Gen. error & Train Acc & Test Acc & Gen. error \\
% \hline
% MLP (no can.)      & 0.889 & 0.461 & 0.428 & 1 & 0.433 & 0.555 \\
% MLP (sort)    & 1.0    & 0.692 & 0.308  & 1 & 0.700 & 0.294 \\
% MLP (hilbert) & 1.0    & \textbf{0.833} & \textbf{0.167}  & 1 & \textbf{0.748} & \textbf{0.252} \\
% \hline
% \end{tabular}
% \end{table}

We next conduct some experiments to corroborate conclusion 1, which claims that, in terms of generalization, group averaging is preferable to canonizations, which are in turn preferable to ignoring invariance altogether.

\textbf{Group averaging on PCA frames} Our first experiment is still with the ModelNet40 classification task, but while the previous experiments investigates canonization w.r.t. permutations, and implicitly relies on the fact that ModelNet models are aligned w.r.t rotations, in this experiment we use the invariant Deepsets \cite{zaheer2017deep} to deal with permutation invariant, and will focus on rotation invariance: we will assume the ModelNet models rotation degree of freedom is initially unknown, and is fixed using the PCA axes as suggested in \cite{punyframe}. Assuming the point cloud singular values are pairwise distinct, this procedure determines the point cloud orientation uniquely up to a $\{-1,1\}^3$ symmetry.

We consider four ways to deal with this symmetry: (i) Pure PCA, where the symmetry is ignored (ii) Skewness canonization, where the global sign of the $x,y$ and $z$ coordinates is chosen so that their third moment is positive (iii) group averaging and (iv) group augmentation where in each epoch a random group element is applied to the input. This last example can be seen as an approximation of group averaging. The results of this experiment, detailed in Table \ref{tab:pca_canonization} and visualized in Figure \ref{fig:loss_pca}, demonstrate a clear performance hierarchy. Group averaging achieves the highest overall accuracy, while group augmentation and skewness canonization yield comparable final results, with all three methods significantly outperforming the pure PCA baseline. Furthermore, as highlighted in Figure \ref{fig:loss_pca}, the group augmentation approach suffers from a noticeably slower convergence rate compared to deterministic skewness canonization. 


% \begin{table}[ht]
% \centering
% \caption{Test accuracy on ModelNet40 with a DeepSets model with orientation determined by PCA. We compare different methods of handling sign ambiguity via averaging, sign augmentation, canonization, or Pure PCA (ignoring the ambiguity).}
% \label{tab:pca_canonization}
% \begin{tabular}{lc}
% \toprule
% \textbf{Method} & \textbf{Accuracy} $\uparrow$ \\
% \midrule
% Sign Averaging  & \textbf{0.760} $\pm$ 0.0025 \\
% Sign Augmentation & 0.742  $\pm$ 0.0022 \\
% Skewness Canonization  & 0.742 $\pm$ 0.0050 \\
% Pure PCA  & 0.729 $\pm$ 0.0062 \\
% \bottomrule
% \end{tabular}
% \end{table}


%\nd{DeepSets on Modelnet10      train: 0.983  test: 0.912}
\textbf{Group averaging for image classification} Lastly, we check our approach on image classification. We take the rotated MNIST \cite{larochelle2007empirical} dataset, and consider the action of the $p4$ group, which is the four-element group generated by $90$ degree rotations. We consider 4 approaches: simple CNN, the canonization from \cite{kaba2023equivariance} with random or (CN(p4 frozen)-CNN) or learned CN(p4)-CNN weights, and a group average over the $4$ elements of the group. As shown in Table \ref{tab:mnist}, the canonizations outperform the simple CNN, but are outperformed by group averaging, as predicted in Conclusion 1. We also see in the table that the accuracy is correlated with the covering numbers of the different methods.
% \begin{table}[h]
% \centering
% \caption{\it Accuracy on Rotated-MNIST using a non-invariant CNN, two canonizations from \cite{kaba2023equivariance}, and group averaging.}
% \label{tab:mnist}
% \begin{tabular}{lcc}
% \toprule
% \textbf{Model} & \textbf{Test Acc.} $\uparrow$ & \textbf{Mean Coverage} $\downarrow$ \\
% \midrule
% CNN  & $94.8 \pm 0.13$ & 0.185  \\
% %SortCNN &  $95.87 \pm 0.0008$\\
% %CN(p4)-CNN(paper) & $97.6 \pm 0.01$ \\
% CN(p4)-CNN & $96.47 \pm 0.24$  & 0.17652 \\
% CN(p4 frozen)-CNN & $95.3 \pm 0.2$ & 0.1765 \\
% AvgCNN & $\textbf{97.3} \pm 0.001$ & \textbf{0.1673}  \\
% \bottomrule
% \end{tabular}
% \end{table}
 
 %
%Regular distance & 0.185   \\
%SortCNN distance    & 0.1718  \\
%CN(p4)-CNN(no train) & 0.17652\\
%CN(p4)-CNN(frozen) & ?\\
%CN(p4)-CNN & 0.1755\\
%Group distance   & 0.1673\\


\begin{table}[h]
    \centering
    \setlength{\tabcolsep}{3pt}
    % --- FIRST TABLE ---
    \begin{minipage}[t]{0.48\textwidth} % [t] aligns them by the top
        \centering

\caption{\it Test accuracy on ModelNet40 with a DeepSets model with orientation determined by PCA, and sign ambiguity handled via averaging, sign augmentation, canonization, or Pure PCA (ignoring the ambiguity).}
\label{tab:pca_canonization}
\begin{tabular}{lc}
\toprule
\textbf{Method} & \textbf{Accuracy} $\uparrow$ \\
\midrule
Sign Averaging  & \textbf{0.760} $\pm$ 0.0025 \\
Sign Augmentation & 0.742  $\pm$ 0.0022 \\
Skewness Canonization  & 0.742 $\pm$ 0.0050 \\
Pure PCA  & 0.729 $\pm$ 0.0062 \\
\bottomrule
\end{tabular}
    \end{minipage}
  \hfill % Pushes the tables to the left and right edges
    % --- SECOND TABLE ---
    \begin{minipage}[t]{0.48\textwidth}
        \centering
       
    \caption{\it Accuracy and mean coverage on Rotated-MNIST which in invariant with respect to the group $p4$ of $90$ rotations. We consider a non-invariant CNN, two canonizations from \cite{kaba2023equivariance}, and group averaging.}
\label{tab:mnist}
\begin{tabular}{lcc}
\toprule
\textbf{Model} & \textbf{Test Acc.} $\uparrow$ & \textbf{Coverage} $\downarrow$ \\
\midrule
CNN  & $94.8 \pm 0.13$ & 0.185  \\
%SortCNN &  $95.87 \pm 0.0008$\\
%CN(p4)-CNN(paper) & $97.6 \pm 0.01$ \\
CN(p4)-CNN & $96.47 \pm 0.24$  & 0.176 \\ %0.17652 \\
CN(p4 frozen)-CNN & $95.3 \pm 0.2$ & 0.176 \\% 0.1765 \\
AvgCNN & $\textbf{97.3} \pm 0.001$ & \textbf{0.167}  \\  %\textbf{0.1673}  \\
\bottomrule
\end{tabular}
    \end{minipage} 
\end{table}

 
\section{Conclusion}
In this paper, we establish a methodology for comparing different canonizations with group averaging and other equivariant models by studying their covering number. We show the covering number of canonizations is at best the covering number of the quotient space (for isometric canonizations), and at worst the covering number of the original domain (for certain discontinuous canonizations). These results suggest that group-averaged models are preferable when computationally feasible. We also show that the covering number of Hilbert canonizations for point clouds is substantially better than that of lexsort canonizations. 

\paragraph{Limitations and Future Work}
A limitation of our analysis is that it relies on upper bounds on the generalization error (Proposition \ref{propGen}), which may not be tight. Nonetheless, our experiments indicate that lower covering numbers do correlate with better generalization and learning. Looking towards future developments, a central contribution of this work is the demonstration that a rigorous theoretical comparison between competing canonization methods can be established via their covering numbers. This framework complements the traditional focus on a model's ability to distinguish between complex symmetric examples, instead highlighting properties that arguably exert a more significant influence on the ultimate success of the learning process. In future work, we intend to apply our framework to other canonizations and extend it to enable a rigorous treatment of the generalization of frame averaging. 

\newpage
\bibliographystyle{plain}
\bibliography{can}
\newpage
\appendix


% \begin{table}[t]
% \centering
% \caption{TBA}
% \label{tab:distances_both}
% \begin{tabular}{lcccc}
% \toprule
% \textbf{Method} & \textbf{RotatedMNIST} \\
% \midrule
% Regular distance & 0.185   \\
% %SortCNN distance    & 0.1718  \\
% %CN(p4)-CNN(no train) & 0.17652\\
% CN(p4)-CNN(frozen) & ?\\
% CN(p4)-CNN & 0.1755\\
% Group distance   & 0.1673\\
% \bottomrule
% \end{tabular}
% \end{table}


\section{Proofs}
In this section, we provide the proofs that are missing in the main text. 


\subsection{Generalization and Covering Numbers}\label{app:gen}
Our generalization results are based on the results of \cite{xu_robustness_2012}: a combination of Theorem 1 and Example 4 from  \cite{xu_robustness_2012} yields the following theorem:  
% Replicating Theorem 1 from 
\begin{theorem}[ \cite{xu_robustness_2012}]\label{thm:xu}
If $\mathcal{Z}$ is compact w.r.t. metric $\rho$, and $l(h_S, \cdot)$ is Lipschitz continuous with Lipschitz constant $c(s)$, i.e.,
\begin{equation*}
|l(h_S, z_1) - l(h_S, z_2)| \le c(S) \rho(z_1, z_2), \quad \forall z_1, z_2 \in \mathcal{Z},
\end{equation*}
 and the training sample set $S$ is generated by $n$ IID draws from $\mu$, then for any $\delta,\gamma > 0$, with probability at least $1 - \delta$ we have
\begin{equation}\label{eq:intermediate}
|\mathcal{L}_{exp}(h_S) - \Lemp(h_S)| \le c(S)\gamma + M \sqrt{\frac{2\mathcal{N}( \mathcal{Z}, \rho,\gamma/2) \ln 2 + 2 \ln(1/\delta)}{n}}.
\end{equation}
\end{theorem}

We now explain how Theorem \ref{thm:MainGeneralizationBound} from the main text follows from Theorem \ref{thm:xu}. We recall Theorem \ref{thm:MainGeneralizationBound}:
\MainGeneralizationBound*
\begin{proof}
We will apply Theorem \ref{thm:xu}, where to apply it to
 to the setting discussed in Theorem \ref{thm:MainGeneralizationBound}, we set
$$z=(x,f(x))  \quad \mathcal{Z}=\{(x,f(x))| \quad x\in X  \} , \quad \rho(z,z')=\rho_X(x,x')+\rho_Y(y,y'), $$
and 
$$l(h_S,z)=\ell(h_S(x),f(x)) . $$
The measure $\mu$ on $X$ is pushed forward to a measure $\hat \mu$ on $\mathcal{Z}$ via the mapping $x\mapsto (x,f(x)) $.

We note that $l$ is Lipscitz with a Lipschitz constant of $c_\ell \max\{1,c_{h_S} \}$ as the following shows: 
\begin{align*}
|l(h_S,z)-l(h_S,z')|&=|\ell(h_S(x),y)-\ell(h_S(x'),y')|\\
&\leq c_\ell \left( \rho_Y(h_S(x),h_S(x'))+\rho_Y(y,y')  \right)\\
&\leq c_\ell \left( c_{h_S}\rho_X(x,x')+\rho_Y(y,y')  \right)\\
&\leq c_\ell \max\{1,c_{h_S} \}\left( \rho_X(x,x')+\rho_Y(y,y') \right)\\
&\leq c_\ell \max\{1,c_{h_S} \} \rho(z,z')
\end{align*}

Next, we note that under the assumptions of Theorem \ref{thm:MainGeneralizationBound},  the covering number of $\mathcal{Z}$ is related to the covering number of $X $ via 
\begin{equation}\label{eq:cover2cover}
\forall \epsilon >0, \quad \N\left( \mathcal{Z},\rho, (1+c_f) \epsilon\right)\leq \N(X,\rho_X,\epsilon) \end{equation}
since we have the upper bound 
\begin{align*}
\rho\left( \left(x,f(x) \right), \left(x',f(x') \right) \right)=\rho_X(x,x')+\rho_Y(f(x),f(x'))
\leq \rho_X(x,x')+c_f\rho_X(x,x')
\end{align*}
and therefore any  cover of $X$ by balls of radius $\epsilon$ is a cover of $\mathcal{Z}$ by balls of radius $(1+c_f)\epsilon $. 

Now to conclude the proof, choose any $\epsilon>0$ and set $\gamma=2(1+c_f)\epsilon $. Since $l$ is $c_\ell \max\{1,c_{h_S} \} $ Lipschitz, we have 
\begin{align*}
|\mathcal{L}_{exp}(h_S) - \Lemp(h_S)|
&\stackrel{\eqref{eq:intermediate}}{\le} c_\ell \max\{1,c_{h_S} \} \gamma + M \sqrt{\frac{2\mathcal{N}( \mathcal{Z}, \rho, \gamma/2) \ln 2 + 2 \ln(1/\delta)}{n}}\\
&=2c_\ell \max\{1,c_{h_S} \} (1+c_f)\epsilon+M \sqrt{\frac{2\mathcal{N}(\mathcal{Z}, \rho, (1+c_f)\epsilon) \ln 2 + 2 \ln(1/\delta)}{n}}\\
&\stackrel{\eqref{eq:cover2cover}}{\le}
2c_\ell \max\{1,c_{h_S} \} (1+c_f)\epsilon+M \sqrt{\frac{2\mathcal{N}( X, \rho_X,\epsilon) \ln 2 + 2 \ln(1/\delta)}{n}}
\end{align*}
which proves Theorem \ref{thm:MainGeneralizationBound}.
\end{proof}

Next, we prove proposition \ref{propGen}.
\generalizationbound*
\begin{proof}
The case where $h_S$ is Lipschitz (but not invariant) follows directly from Theorem \ref{thm:MainGeneralizationBound}.

Next, we  consider the case that $h_S=\tilde{h}_S\circ c$ where $\tilde{h}_S$ is $c_{h_S}$ lipschiz. We want to consider the domain to be $c(K)$, so consider the following distribution measure on $c(K)$: 
    \begin{align*}
        \tilde{\mu}(A)= \mu(c^{-1}(A))
    \end{align*}
    Note that $$\forall k \in c(K), h_s(k)=\tilde{h}_s(k) $$
    Next, sample $n$ IID samples from $K$ denoted by $x_1,x_2,\ldots,x_n$ and note that $c(x_1),c(x_2),\ldots,c(x_n)$ are an IID samples from $c(K)$ by $\tilde{\mu}$ and denote them by $s_1,s_2,\ldots,s_n$. 
    Note that 
    \begin{align*}
        \Lemp(h_S)=\frac{1}{n}\sum_{i=1}^n \ell(h_S(x_i),f(x_i)) = \frac{1}{n}\sum_{i=1}^n \ell(\hat{h}_S(s_i),f(s_i)) = \Lemp(\hat{h}_S)
    \end{align*}
    So, the empirical losses are the same, and next we show the expected losses: 
    \begin{align*}
        \mathcal{L}_{exp}(\tilde{h}_S) &= \mathbb{E}_{x \sim \tilde{\mu}} \ell(\tilde{h}_S(x),f(x)) = \int_{c(K)} \ell(\tilde{h}_S(x),f(x)) d\tilde{\mu} = \int_{c^{-1}(c(K))} \ell(\tilde{h}_S(c(x)),f(c(x)) d\mu  \\ 
        &=\int_{K} \ell(h_S(x)),f(x) d\mu = \mathcal{L}_{exp}(h_S)
    \end{align*}
    Finally, we apply Theorem \ref{thm:MainGeneralizationBound} on $\tilde{h}_S$ and as it's $c_{h_S}$ lipschiz we obtain that with probablity at least $\delta$ holds that: 
    \begin{align*}
        |\mathcal{L}_{exp}(\tilde{h}_S) - \Lemp(\tilde{h}_S)| \le 2c_\ell \max\{1,c_{h_S} \} (1+c_f)\epsilon+M \sqrt{\frac{2\mathcal{N}( c(K), \rho_X,\epsilon)\ln 2 + 2 \ln(1/\delta)}{n}}
    \end{align*}
    And as $|\mathcal{L}_{exp}(\tilde{h}_S) - \Lemp(\tilde{h}_S)| = |\mathcal{L}_{exp}(h_S) - \Lemp(h_S)|$ 
    
    We obtain that \begin{align*}
        |\mathcal{L}_{exp}(h_S) - \Lemp(h_S)| \le 2c_\ell \max\{1,c_{h_S} \} (1+c_f)\epsilon+M \sqrt{\frac{2\mathcal{N}( c(K), \rho_X,\epsilon)\ln 2 + 2 \ln(1/\delta)}{n}}
    \end{align*}
    In the second case, where $h_S$ is $c_{h_S}$ lipshiz and invariant, we want to consider the domain $K/\G$ and define $\pi: K \rightarrow K/\G$ by $\pi(k)= [k]$. Define the following distribution on $K/\G$ by: 
    \begin{align*}
        \tilde{\mu}(A)= \mu(\pi^{-1}(A))
    \end{align*}
    Define $$\tilde{h}_S,\tilde{f}: k/ \G \rightarrow \R$$ by $$\tilde{h}_S([x])=h_S(x), \tilde{f}([x]) = f(x)$$
    Note that by invariance, those functions are well defined. 
    Next, sample $n$ IID samples from $K$ denoted by $x_1,x_2,\ldots,x_n$ and note that $\pi(x_1),\pi(x_2),\ldots,\pi(x_n)$ are an IID samples from $K/\G$ by $\tilde{\mu}$ and denote them by $s_1,s_2,\ldots,s_n$. 
    Note that 
    \begin{align*}
        \Lemp(h_S)=\frac{1}{n}\sum_{i=1}^n \ell(h_S(x_i),f(x_i)) = \frac{1}{n}\sum_{i=1}^n \ell(\hat{h}_S(s_i),\tilde{f}(s_i)) = \Lemp(\hat{h}_S)
    \end{align*}
    So, the empirical losses are the same, and next we show the expected losses: 
    \begin{align*}
        \mathcal{L}_{exp}(\tilde{h}_S) = \mathbb{E}_{x \sim \tilde{\mu}} \ell(\tilde{h}_S(x),\tilde{f}(x)) = \int_{K/\G} \ell(\tilde{h}_S(x),\tilde{f}(x)) d\tilde{\mu} = \int_{\pi^{-1}(K/\G)} \ell(\tilde{h}_S(\pi(x)),\tilde{f}(\pi(x))) d\mu = \\
        \int_{K} \ell(\tilde{h}_S([x]),\tilde{f}([x])) d\mu
        = \\
        \int_{K} \ell(h_S(x)),f(x) d\mu = \mathcal{L}_{exp}(h_S)
    \end{align*}
    Note that $\tilde{f}$ is $c_f$ lipshiz function and $\tilde{h}_S$ is $c_{h_S}$ lipshiz function. 
    Finally, we apply Theorem \ref{thm:MainGeneralizationBound} on $\tilde{h}_S$ and $\tilde{f}$ and as it's $c_{h_S}$ lipschiz we obtain that with probablity at least $\delta$ holds that: 
    \begin{align*}
        |\mathcal{L}_{exp}(\tilde{h}_S) - \Lemp(\tilde{h}_S)| \le 2c_\ell \max\{1,c_{h_S} \} (1+c_f)\epsilon+M \sqrt{\frac{2\mathcal{N}(K/ \G, \rho_\G,\epsilon)\ln 2 + 2 \ln(1/\delta)}{n}}
    \end{align*}
    And as $|\mathcal{L}_{exp}(\tilde{h}_S) - \Lemp(\tilde{h}_S)| = |\mathcal{L}_{exp}(h_S) - \Lemp(h_S)|$ 
    
    We obtain that \begin{align*}
        |\mathcal{L}_{exp}(h_S) - \Lemp(h_S)| \le 2c_\ell \max\{1,c_{h_S} \} (1+c_f)\epsilon+M \sqrt{\frac{2\mathcal{N}(k/\G, \rho_\G,\epsilon)\ln 2 + 2 \ln(1/\delta)}{n}}
    \end{align*}
\end{proof}
\subsection{Proof for the equivariant case}
\begin{restatable}{proposition}{Equigeneralizationbound}
Let $(K,\rho,\G)$ be a module. Let  $(Y,\rho_Y) $ be a metric space, and assume $K$ is compact. Let $f:K\to Y $ be a $c_f$  Lipschitz function \textbf{which is $\G $ equivariant}, and let $\ell:Y\times Y \to [0,M] $ be a $c_\ell $ Lipschitz function and $\G$ invariant, then for any $\epsilon,\delta>0$ and natural $n$,  for a training sample set $S$ generated by $n$ IID draws from the distribution $\mu$,  with probability of at least $ 1-\delta$, we have 
$$|\mathcal{L}_{exp}(h_S) - \Lemp(h_S)| \le 2c_\ell \max\{1,c_{h_S} \} (1+c_f)\epsilon+M \sqrt{\frac{2\textcolor{blue}{\N(h_S)} \ln 2 + 2 \ln(1/\delta)}{n}}$$

where 
\begin{equation}
\textcolor{blue}{\N(h_S)}=\threepartdef{\N(K,\rho,\epsilon)}{h_s \text{ is } c_{h_s} \text{ Lipschitz }}{\N(c(K),\rho,\epsilon)}{h_s=\tilde{h}_s\circ c, \tilde{h}_s \text { is } c_{h_s} \text{ Lipschitz, and }  c \text{ is a canon.}} {\N(K/\G,\rho_{\G},\epsilon)}{h_s \text{ is } c_{h_s} \text{ Lipschitz and  } \G \text{ Equivariant }} 
\end{equation}
\begin{proof}
    We begin with the second case. Define $l(h_S,z)=\ell(h_S(x),f(x))$. We frist show it's $\G$ invariant w.r.t $z$:
    \begin{align*}
        l(h_S,g\cdot z)=\ell(h_S(g\cdot x),f(g\cdot x)) = \ell(g\cdot h_S( x),f(g\cdot x))
    \end{align*}
\end{proof}
\end{restatable}



\subsection{Proofs for Section \ref{sec:evaluation}}
We restate and prove Proposition \ref{prop:cardinality}.
\cardinality*
\begin{proof}
\label{upper_bound}
Let $C$ be an $\epsilon$-cover of $K/\G$ under $\rho_{\G}$, and define
\[
\widehat{C}:=\{g\cdot y \mid y\in C,\ g\in \G\}.
\]
Then of course $|\widehat{C}|\le |\G|\cdot |C|$.

We claim that $\widehat{C}$ is an $\epsilon$-cover of $K$ under $d$.
Let $x\in K$. Since $C$ covers $K/\G$, there exists $y\in C$ such that
\[
\rho_{\G}([x],[y])\le \epsilon.
\]
Because $\G$ is finite, the minimum in the definition of $\rho_{\G}$ is attained, so there exists $g\in \G$ such that
\[
\rho(x, g\cdot y)\le \epsilon.
\]
By construction, $g\cdot y\in \widehat{C}$, hence $\widehat{C}$ covers $K$.
Therefore $\N(K,\rho,\epsilon)\le |\widehat{C}|\le |\G|\cdot |C| = |\G|\cdot \N(K/\G,\rho_{\G},\epsilon)$.
\end{proof}
We next prove Claim \ref{claim:isometry}.
\isometry*
\label{Proof_isometry}
\begin{proof}
    As it always holds that \[\N(K/\G,\rho_{\G},\epsilon)\leq \N(c(K),\rho,\epsilon)\] we just need to prove the opposite direction.
    
    Let $C$ be a cover of $K/\G$ with $\N(K/\G,\rho,\epsilon)$ elements.
    Consider \[\hat{C}:=\{c(x)|x \in C \}\] and we will prove that $\hat{C}$ is an $\epsilon$ cover of $c(K)$.
    
    Let $y \in c(K)$, then by definition $\exists x \in K: y = c(x)$. By the definition of cover $\exists z\in C$ such that $\rho([x],[z])\leq \epsilon$. But as $\rho_{\G}([x],[z])=\rho(c(x),c(z))$ we obtain that $\rho(y,c(z))\leq \epsilon$. So $\hat{C}$ is an $\epsilon$ cover of $c(K)$ as desired. 
\end{proof}

\isometries*
\label{Proof_isometries}
\begin{proof}
The  canonization $c(x)=|x|$ is an isometry as
$$\forall x,y \in K, \rho_\G([x],[y]) = \min\{|x-y|,|x+y|\} = ||x|-|y|| = |c(x)-c(y)|$$

The fact that sorting is an isometry with respect to the metric $\rho_\G$ induced from a $p$ norm (which is the $p$-Wasserstein metric) is shown in 
 Lemma 4.2 in  \cite{sortLemma}. 

consider the action of $\RR^d$ on $\RR^{d\times n}$ by translating all columns of a matrix in $\RR^{d\times n}$ by a vector in $\RR^d$. Let $\rho$ be the metric induced from the Frobenius norm $\rho(X,Y)=\|X-Y\|_F $, and let $c$ be the centralization mapping $c(X)=X-\frac{1}{n}X1_{n\times n} $. This mapping is a canonization and an isometry. To see it is an isometry, note that we have for any canonization $\rho(c(X),c(Y))\leq \rho_\G\left([X],[Y]\right) $, and we get the converse inequality by noting that $c$ is the orthogonal projection from the inner product space $\RR^{d\times n} $ onto the subspace $\{X\in \RR^{d\times n}| \quad X1_{n\times n}=0 \} $. Thus, let $g\in \G$ be a translation vector such that $\rho_\G([X],[Y])=\rho(X,gY) $. Then 
$$\rho_\G([X],[Y])=\rho(X,gY)=\|X-gY\|_F\leq \|c(X)-c(gY)\|_F=\|c(X)-c(Y)\|_F=\rho(c(X),c(Y)) $$

\end{proof}


We restate and prove \ref{prop:poor}.
\poorcanonization*
\begin{proof}
\label{poor_proof}
$L$ is a partial limit of $c$ if there is a sequence $x_n \rightarrow x$ such that $c(x_n)\rightarrow L$. We claim that $L$ must be in the orbit of $x$. Indeed, $c(x_n)=g_nx_n $ for appropriate $g_n \in G$. Since $\G$ is finite, by passing to a subsequence, we can assume that $g_n=g\in G$ does not depend on $n$. Since $x\mapsto gx$ is an isometry, we know that $gx_n \rightarrow gx=L $.  In particular, $gx\neq x $ for every $g\neq e$. Choose some $\epsilon < min_{g \neq e} |x - g\cdot x|$, and define the $G$ invariant set $K=K(\epsilon)$ by $K := \cup_{g \in G} B(g\cdot x, \epsilon)$. Note that $c(K)$ contains all partial limits of $c$ at $x$, and therefore $gx\in c(K) $ for all $g\in G$. Since the distance between any two elements in the orbit of  $ x$ is larger than $\epsilon$, an $\epsilon$ cover of $c(K)$ must contain at least $|\G| $ elements. On the other hand, $K$ itself is a union of $|\G|$ balls of radius $\epsilon$, and therefore
$$\N(c(K),\rho,\epsilon)\leq \N(K,\rho,\epsilon)\leq |\G| \leq \N(c(K),\rho,\epsilon) , $$
and therefore all these inequalities are equalities and the covering number is exactly $|\G|$.
   
Finally, note that $\rho(K/G,\rho_G,\epsilon)=1$ as one ball of radius $\epsilon$ is enough for covering the quotient space, as all $|\G|$ balls are the same up to group action.
\end{proof}
\subsection{Proofs for Section \ref{sec:case_study}}
For proving Lemma \ref{lem:quotient_covering}, we need the following lemma: 
\begin{lemma}
    \label{comb_proof}
    Let $S := [m]^n$ all vectors with $n$ elements, each between $1$ and $m$. Define two vectors $v\equiv u \iff \exists \pi \in S_n : u = \pi \cdot v$. Then the number of equivalent classes is exactly $\binom{n+m-1}{n}$.
\end{lemma}
\begin{proof}
    Define $\forall i \in [m]$, $c^u_i$ to be the number of times $i$ appears in $u$. 
    \newline
    Note that $u \equiv v \iff \forall i\in [m]: c^u_i=c^v_i$. 
    So we want to ask how many solutions there are to the following system:
    \begin{align*}
        \sum^{m}_{i=1}c_i = n \\
        \forall i,c_i \geq 0
    \end{align*}
    By \cite{stanley2011enumerative} we know the number of solutions to the equation is $\binom{n+m-1}{n}$. 
\end{proof}
We now restate and prove \ref{lem:quotient_covering}.
\quotientcovering*
\begin{proof}
Set $k=\lceil 1/(2\epsilon) \rceil $. The set $K$ can be covered by $k^d$ hypercubes of diameter $k$ and radius $1/(2k)\leq \epsilon $. Let $C$ denote the collection of centers of these hypercubes, namely 
\begin{align*}
    C=\left\{\left[\frac{2i_1+1}{2k} ,\frac{2  i_2+1}{2k} ,\ldots,\frac{2 i_n+1}{2k} \right]| \quad  i_j \in \{0,1,\ldots,k-1 \},1\leq j \leq n \right\}   .
\end{align*}    
Applying the quotient map  $\pi(c) = [c]$ we obtain points $\pi(C)$ in the quotient space which are an $\epsilon$ cover of $K/\G$. Note that points in $C$ which are permutations of each other are mapped by $\pi$ to the same orbit. Therefore, the cardinality of the cover $\pi(C)$ is the number of equivalence classes with respect to the relation on $C$ defined by the group:  $c_1 \equiv c_2 \iff \exists \pi \in S_n : c_1 = \pi(c_2)$. Using Lemma \ref{comb_proof}, we know that the number of equivalence classes is \begin{equation*}
|\pi(C)|=    \binom{n+k^d-1}{n}
\end{equation*} 
\end{proof}
We now restate and prove \ref{lem:lexsortcovering}.
\lexsortcovering*
\begin{proof}
    \label{lexsortcovering}
    We consider the set $U \subset [0,1]^{ d \times n}$ defined by 
\begin{align*}
    U= \{X\in [0,1]^{d\times n}| \quad X_{11}=X_{12}=\ldots=X_{1n}\}
\end{align*}
The set $U$ is not a subset of $\csort(K)$, but it is a subset of its closure $\overline{\csort(K)}$. This is because for every  $X \in U$ we can apply an arbitrarily small perturbation of the first row of $X$ so that it is ordered from small to large, and the $X$ will be unaffected by the canonization. We then obtain
$$ \N(\csort(K),\rho,\epsilon)=\N(\overline{\csort(K)},\rho,\epsilon)\geq \N(U,\rho,\epsilon)\stackrel{(*)}{=}\left(\frac{1}{2\epsilon}\right)^{(d-1)\cdot n + 1},$$
where $(*)$ follows from the fact that the hypercube $[0,1]^{(d-1)\cdot n + 1} $ can be mapped isometrically (in the $\infty$ norm) to $U$ via $f:[0,1]^{(d-1)\cdot n + 1} \rightarrow U$ by repeating the first entry $n-1$ times, obtaining an $n \cdot d$ vector, and then transposing to be of shape $n \times d$, and from the fact that the covering number of the hypercube is  
$\left(\frac{1}{2\epsilon}\right)^{(d-1)\cdot n +1}$.
\end{proof}
\subsection{Hilbert Canonization}\label{app:hilbert}
In this appendix, we give a formal definition of Hilbert curves and a full proof of Theorem \ref {thm:hilbert}.

The Hilbert curve $H(x)$ is defined as a limit of the function $H_m:[0,1]\to [0,1]^d$, and these $H_m$ are our main focus. We will now define these curves, loosely following the terminology from 
\cite{Hilbert}. Firstly, for a given integer $L $, we consider a partition of the unit interval into $2^L$ disjoint intervals: the semi-closed intervals $I_L(\ell)=[\ell 2^{-L}, (\ell+1)2^{-L})$ for $\ell=0,\ldots,2^{L}-2$, and the closed interval $I(2^L-1)=[(2^L-1)2^{-L},1] $. We denote the collection of all these intervals by $\mathcal{I}_L $. For a natural $d>1$, we also consider the partition of the unit hypercube $[0,1]^d$ into $2^{L\cdot d} $ disjoint hypercubes of the form 
$E_L^d(\vec \ell)=I(\ell_1)\times \ldots \times I(\ell_d)$. We denote the collection of these hypercubes by $\mathcal{E}_L^d $. 

Next, for any given $m$, we define  $H_m$ to be a  bijection from $\mathcal{I}_{d\cdot m} $ to $\mathcal{E}_{m}^d $, such that (i) adjacent intervals are mapped to adjacent hypercubes and (ii) if $I_{d\cdot (m+1)}(k')\subseteq I_{d\cdot m}(k) $ then 
$H_{m+1}\left(I_{d\cdot (m+1)}(k')\right)\subseteq H_m\left(I_{d\cdot m}(k) \right) $ . An illustration of $H_m, m=1,2,3$ for $d=2$ is given in Figure \ref{fig:hilbert}. Each $H_m$ maps a finite collection of intervals bijectively to a finite collection of hypercubes. This naturally induces a bijection from the grid defined by the intervals centroids, to the grid defined by the hypercube centroids. We denote these grids by $\mathbb{A}[d\cdot m]\subseteq [0,1] $ and $\mathbb{B}[d,m]\subseteq \RR^d $, respectively.

The Hilbert curve $H$ is obtained from $H_m$ by a limiting procedure and shown to be Holder continuous \cite{Hilbert}. Following similar ideas we show that each Holder function $H_m$ is Holder continuous with constant that do not depend on $m$. This property will be crucial for the analysis of the covering number of the Hilbert canonization.
\begin{restatable}{lemma}{Holder}\label{lem:Holder}
For every natural $m,d$, we have that 
$$\forall x,y\in \mathbb{A}[d\cdot m], \quad  \|H_m(x)-H_m(y)\|_\infty \leq 4|x-y|^{1/d} $$
\end{restatable}
\begin{proof}
    Let $x\neq y$ be points in the one dimensional grid $\mathbb{A}[d\cdot m]$, and assume without loss of generality that $x<y$. Note that $y-x\geq \frac{1}{2^{d\cdot m}}$, since this is the minimal distance between points in the grid. On the other hand, $y-x\leq 1$. Therefore, there exists some $k$ with $m-1\leq k \leq 0 $, such that 
    \begin{equation}\label{eq:interval} 
    2^{-d\cdot (k+1)}\leq |x-y|\leq 2^{-d \cdot k}.
    \end{equation}
   It follows that $x$ and $y$ are in two adjacent intervals in the $k$-th level (or in the same interval), and by the nesting property of Hilbert's curve, we know that $H_m(x)$ and $H_m(y)$ are in two adjacent hypercubes (or in the same hypercube). Therefore, their distances is at most the sum of the diameters of those two cubes, so
    \begin{align*}
    \|H_m(x)-H_m(y)\|_\infty\leq 2\cdot 2^{-k}\stackrel{\eqref{eq:interval}}{\leq} 4|x-y|^{1/d}
    \end{align*}
\end{proof}



\textbf{Hilbert Canonization}
As a first step, we define a canonization not on all of $[0,1]^{d\times n}$, but rather on the finite set of $n$ tuples of grid elements $\mathbb{B}^n[d,m] $. This is done by mapping each tuple-element to the unit interval using $H_m^{-1} $, then sorting the resulting $n$-dimensional array, and mapping back to the $d$-dimensional grid using $H_m$. We denote this canonization by $c_m$. We can write $c_m=H_m \circ \mathrm{sort}\circ H_m^{-1}$, with the understanding that $H_m$ and  $H_m^{-1} $ are applied elementwise to $n$-tuples of grid elements. 

To extend $H_m$ to all of $K=[0,1]^{d\times n}$, we employ the following procedure: For a given $X=(x_1,\ldots,x_n)\in K  $, (i) \textbf{Rounding:} Each $x_j$ resides in a unique hypercube in $ \mathcal{E}_m^d$. Replace $x_j$ with the centroid of this cube $y_j\in \mathbb{B}^n[d,m] $. (ii) \textbf{Sorting:} find the permutation $\tau \in S_n$ which sorts $\left(H_m^{-1}(y_1),\ldots,H_m^{-1}(y_n) \right)$ from small to large. (iii) \textbf{Canonizing} Apply the permutation $\tau$ to $X$. We note that if $X$ contains several points in the same hypercube, the permutation $\tau$ is not uniquely defined. In this case, the order of points with the hypercube will be determined by lexicographical sorting. 

We now restate and prove our theorem regarding the covering number of the Hilbert canonization:


\hilbert*
\begin{proof}
We first begin by a reduction from covering $K$ to covering the grid $\mathbb{B}^n[d,m] $ For convenience we abbreviate $\mathbb{B}^n:=\mathbb{B}^n[d,m] $ . We note that for any  $X \in K$, the matrix $Y$ obtained by the rounding procedure satisfies $\|X-Y\|_\infty\leq 2^{-m-1} $. This stays true also after applying the canonization,  $\|c_m(X)-c_m(Y)\|_\infty\leq 2^{-m-1} $, since the canonization will apply the same permutation to both $X$ and $Y$. It follows that 
$$\N\left(c_m\left(K \right),\rho_\infty,\epsilon \right)\leq \N\left(c_m\left(\mathbb{B}^n\right),\rho_\infty,\epsilon-2^{-m-1} \right) .$$
It remains to bound the right-hand side:
\begin{align*}
 \N\left(c_m\left(\mathbb{B}^n\right),\rho_\infty,\epsilon-2^{-m-1} \right)&=\N\left(H_m \circ \mathrm{sort} \circ H_m^{-1}\left(\mathbb{B}^n\right),\rho_\infty,\epsilon-2^{-m-1} \right) \\
 &\stackrel{(*)}{\leq} \N\left( \mathrm{sort} \circ H_m^{-1}\left(\mathbb{B}^n\right),\rho_\infty,\delta \right)\\
 &\leq \N\left( \mathrm{sort} \left( [0,1]^n \right),\rho_\infty,\delta \right)\\
 &\stackrel{(**)}{\leq}  \binom{n+\lceil 1/ (2\delta) \rceil-1}{n}
\end{align*}
Where (*) follows from the Holder properties of $H_m$ proven in
Lemma \ref{lem:Holder}, and (**) follows from the fact that one-dimensional sorting is an isometry (Proposition \ref{prop:isometries}) and the upper bound on the covering number of $[0,1]^n/\G$ (Lemma \ref{lem:quotient_covering} ).
\end{proof}
\begin{lemma}
    \label{covering_cube}
    Be $K = [0,1]^{n \times d}$ and $\epsilon=\frac{1}{2\cdot k}>0$ such that $k\in \bbN$. Then $\N(K,d_{\infty},\epsilon)=\frac{1}{(2\epsilon)^{nd}}$
\end{lemma}
\begin{proof}
    It's easy to see that each line we can cover using $\frac{1}{2\cdot \epsilon}$ $\epsilon$ sized lines, so using the product of size covering for $n \cdot d$ times we get a covering for $K$ that uses $\frac{1}{(2\epsilon)^{nd}}$ balls.
    On the other hand, denote by $F$ some other covering. 
    Note that $K \subseteq \cup_{c \in F} B(c, \epsilon)$
   So $1=\mu(K)\leq \mu(\cup_{c \in F} B(c, \epsilon))\leq \sum_{c \in F} \mu(B(c, \epsilon))=|F|\cdot (2\epsilon)^{n \cdot d}$ so $\frac{1}{(2\cdot \epsilon)^{nd}}\leq |F|$. 
   Thus, we got the lower bound and the upper bound as desired.
\end{proof}
\section{Additional Experiments}
\label{app:additional_exp}

\subsection{Laplacian Eigenvector Sign Ambiguity: Augmentation vs.\ Canonization on \texttt{ogbg-molpcba}}
\label{subsec:lpe_molpcba}

This appendix presents an additional empirical comparison between deterministic
canonization and training-time random group augmentation, on a setting
distinct from the main-text experiments and complementary to the PCA
experiment of Section~\ref{sec:case_study}. The goal is to ask whether
the empirical behaviour of these two strategies on a standard graph
property prediction benchmark is qualitatively consistent with our
Proposition~\ref{propGen} and the covering-number inequality~\eqref{eq:coveringIneq};
we do not claim a quantitative test of the inequality, since test AP
on a finite-sample-trained GIN mixes optimisation, capacity and
generalisation effects that the inequality does not isolate. The
experiment is also consistent with the empirical observation cited in
the Related Work that randomized SMILES can outperform canonized SMILES
\cite{ArusPous2019,Bjerrum2017SMILESEA}.

\paragraph{Module.}
We study graph property prediction on \texttt{ogbg-molpcba} using a
$5$-layer Graph Isomorphism Network (GIN) with Laplacian positional
encoding (LapPE). For each graph $G$ with normalized Laplacian
$L_G$, we compute the $k$ smallest non-trivial eigenpairs
$\{(\lambda_i, v_i)\}_{i=1}^{k}$ and concatenate the eigenvectors as a
node-level positional encoding. The resulting LapPE is invariant to
graph isomorphism only up to the eigenvector \emph{ambiguity group}
\[
  \G_{\mathrm{LapPE}} \;=\; \prod_{i\,:\,m_i=1}\{-1,+1\} \,\times\, \prod_{i\,:\,m_i>1} O(m_i),
\]
where $m_i$ denotes the multiplicity of eigenvalue $\lambda_i$: a sign
flip is admissible for each simple eigenvalue and an arbitrary orthogonal
change of basis is admissible within each multiplicity-$m_i$ block. On
molecular graphs the spectrum is almost everywhere simple, so in
practice $\G_{\mathrm{LapPE}}$ reduces to $\{-1,+1\}^{k}$ for almost all
graphs. The natural questions are then (i) whether to canonize this
ambiguity, and (ii) whether augmentation-based approximations of group
averaging behave consistent with our framework.

\paragraph{Three methods compared.}
We compare three strategies for resolving $\G_{\mathrm{LapPE}}$, all
implemented as faithful numpy ports of the reference torch code from
the Laplacian canonization release of \cite{ma2023laplacian}.
The first, \texttt{map} (Maximal Axis Projection), resolves a
simple-eigenvalue sign by forming the rank-one projector
$P_i = u_i u_i^{\top}$, grouping the coordinate indices $j$ by the
rounded column norms $\|P_i[:,j]\|$ ($14$-decimal rounding), and
walking these groups from smallest to largest norm to find the first
group whose indicator vector $X$ (the sum of standard basis vectors
over the group, plus a uniform shift) gives $\|P_i X\|\!\neq\!0$; the
canonical sign is then $\mathrm{sign}(u_i^{\top}X)$. For a
multiplicity-$m_i$ block, \texttt{map} forms $P = U_i U_i^{\top}$,
selects the $m_i$ \emph{largest}-norm coordinate groups, and builds
the canonical orthonormal basis incrementally by projecting each
group-indicator vector onto the residual orthogonal complement
(with plain QR inside the complementary-space step).
The second, \texttt{oap} (Orthogonalized Axis Projection), differs
from \texttt{map} in two places: the sign step rounds the column
norms to $6$ decimals (coarser tie-grouping under near-degenerate
norms), and the basis step replaces the largest-norm groupings of
\texttt{map} by a hash-based grouping of columns and uses signed-QR
in the complementary-space iteration.
The third, \texttt{random\_augmented}, flips the sign of each
eigenvector with probability $1/2$ at every training-time forward
pass, so the network sees a uniformly random element of
$\{-1,+1\}^{k}$ at each step. This is a finite-group augmentation
that approximates group averaging over $\G_{\mathrm{LapPE}}$ without
summing over its $2^{k}$ elements explicitly. At evaluation time the
\texttt{random\_augmented} model is fed the deterministic eigsh-sign
convention from the cache (no fresh random draws), yielding a single
deterministic prediction per graph; this is the number reported in
the \texttt{random\_augmented} column of
Table~\ref{tab:lpe_molpcba_main}. Note that this is an
\emph{adversarial} test-time deployment of an averaging-trained model
(it scores a single point of the orbit rather than the orbit mean);
the \textsc{aug-$K$} column reports the alternative Reynolds-style
sample of $K$ orbit elements, and we discuss the gap between the
two below.

\paragraph{Test-time averaging \textsc{aug-$K$}.}
For \texttt{random\_augmented} models we additionally evaluate a
test-time \emph{Reynolds-style} averaged metric. At evaluation, for
each test graph we draw $K$ independent group elements
$g_1,\ldots,g_K \in \G_{\mathrm{LapPE}}$ (each $g_j$ a sign-flip pattern
on the simple-eigenvalue subspace and a Haar-random orthogonal matrix
on each multiplicity-$m$ block); we forward-pass the model on each
sign-flipped LapPE; we average the predicted per-task probabilities;
and we score the averaged predictions with the official OGB AP
evaluator. We use $K\!=\!8$ at $k\!=\!3$ (which approximately covers
$\{-1,+1\}^{3}$ under uniform i.i.d.\ sampling: $8$ draws cover
$\sim\!65\%$ of the $8$-element group on average) and $K\!=\!16$ at
$k\!\in\!\{8,16\}$ (sparse Monte-Carlo over $\{-1,+1\}^{8}$ and
$\{-1,+1\}^{16}$ respectively). On \texttt{ogbg-molpcba} the
spectrum is almost everywhere simple, so the
multiplicity-block factor $\prod_i O(m_i)$ is degenerate in
practice and the $g_j$ samples reduce to sign-flip patterns.
This is the empirical group-averaging operator of
Section~\ref{sec:gen}, restricted to a finite sample.

\paragraph{Hyperparameters and protocol.}
GIN, $5$ layers, hidden dim $h\!\in\!\{16,128,512\}$, dropout $0.5$,
batch size $32$, optimizer Adam with cosine learning-rate schedule and
initial rate $10^{-3}$, $3$ seeds per cell ($\{0,1,2\}$),
eigval-scaling off, eigenvector cache size $15$. Checkpoint selection:
the model state at the epoch with the highest validation AP across
all training epochs is saved (ties broken by earliest epoch); the
final test AP reported in
Tables~\ref{tab:lpe_molpcba_main}~and~\ref{tab:lpe_molpcba_k16} is
this checkpoint's score on the OGB test split. Software stack:
PyTorch $2.10.0$, CUDA $12.8$, \texttt{torch\_geometric} $2.7.0$,
\texttt{numpy} $2.2.6$, \texttt{scipy} $1.15.3$, \texttt{ogb} $1.3.6$
(OGB evaluator: \texttt{Evaluator(name=`ogbg-molpcba')}). Per-arm
seeding via \texttt{torch.manual\_seed}, \texttt{numpy.random.seed},
and \texttt{random.seed}; \texttt{cudnn.deterministic} is not
enforced, so two runs at the same seed on different hardware are not
guaranteed to be bit-identical (we observe the resulting variability
absorbed into the seed std at the magnitudes shown in
Table~\ref{tab:lpe_molpcba_k16}). Standard deviations reported below
are sample standard deviations ($\mathrm{ddof}\!=\!1$). Code, raw
\texttt{results.json} files, the \texttt{orbit\_stability.json} files
behind Table~\ref{tab:lpe_molpcba_orbit}, and reproduction commands
are released at the project repository.

\paragraph{Compute budget.}
At $k\!\in\!\{3,8\}$ all three arms are trained for $200$ epochs with
early-stop patience $15$. At $k\!=\!16$, where validation AP for
\texttt{random\_augmented} continues to improve well past $200$
epochs, all three arms are trained for $500$ epochs with early
stopping disabled (patience $999$). Tables~\ref{tab:lpe_molpcba_main}
and~\ref{tab:lpe_molpcba_k16} therefore report a matched-compute
comparison at every $(k,h)$ cell. Both schedules are released in the
project repository.

\paragraph{Results.}
Table~\ref{tab:lpe_molpcba_main} reports mean test AP at all nine
$(k,h)$ cells for the three methods, alongside the
\textsc{aug-$K$} test-time-averaged metric for \texttt{random\_augmented}.
Table~\ref{tab:lpe_molpcba_k16} zooms in on the $k\!=\!16$ slice with
sample standard deviations.

\begin{table}[h]
\centering
\caption{Mean test AP on \texttt{ogbg-molpcba} at matched compute,
$3$ seeds per cell ($k\!\in\!\{3,8\}$ at $200$ epochs / patience $15$,
$k\!=\!16$ at $500$ epochs / patience $999$ for all three arms).
Per-cell row-best margins are mostly within seed std and none survives
a Bonferroni correction across the nine cells at $\alpha\!=\!0.05$
(see Table~\ref{tab:lpe_molpcba_k16} for the $k\!=\!16$ Welch
$p$-values); the bolding therefore marks the nominal row-best mean,
not a significance claim. The \textsc{aug-$K$} column reports
test-time averaging over $K$ ambiguity-group draws applied to the
\texttt{random\_augmented}-trained model (see text).}
\label{tab:lpe_molpcba_main}
\begin{tabular}{rrcccc}
\toprule
$k$  &  $h$  & \texttt{map} $\uparrow$ & \texttt{oap} $\uparrow$ & \texttt{random\_augmented} $\uparrow$ & \textsc{aug-$K$} \\
\midrule
 3   &   16  & $0.0626$        & $0.0641$        & $\mathbf{0.0662}$           & $0.0663$ ($K\!=\!8$)  \\
 3   &  128  & $0.1704$        & $0.1682$        & $\mathbf{0.1762}$           & $0.1765$ ($K\!=\!8$)  \\
 3   &  512  & $0.2397$        & $\mathbf{0.2431}$ & $0.2395$                  & $0.2398$ ($K\!=\!8$)  \\
\midrule
 8   &   16  & $0.0622$        & $0.0605$        & $\mathbf{0.0636}$           & $0.0637$ ($K\!=\!16$) \\
 8   &  128  & $0.1705$        & $0.1754$        & $\mathbf{0.1766}$           & $0.1766$ ($K\!=\!16$) \\
 8   &  512  & $0.2344$        & $0.2371$        & $\mathbf{0.2395}$           & $0.2400$ ($K\!=\!16$) \\
\midrule
16   &   16  & $\mathbf{0.0716}$ & $0.0703$      & $0.0691$                    & $0.0693$ ($K\!=\!16$) \\
16   &  128  & $\mathbf{0.1923}$ & $0.1836$      & $0.1916$                    & $0.1921$ ($K\!=\!16$) \\
16   &  512  & $0.2468$        & $0.2419$        & $\mathbf{0.2483}$           & $0.2481$ ($K\!=\!16$) \\
\bottomrule
\end{tabular}
\end{table}

\begin{table}[h]
\centering
\caption{Matched-compute $k\!=\!16$ slice of
Table~\ref{tab:lpe_molpcba_main}: all three arms trained for $500$
epochs with patience $999$. Reported are mean $\pm$ sample standard
deviation over $3$ seeds, plus the two-sided Welch $t$-test $p$-value
for the row-best mean against the row-second-best mean. None of the
three row-best margins reaches $p\!<\!0.05$, and after a Bonferroni
correction across the nine cells in
Table~\ref{tab:lpe_molpcba_main} ($\alpha\!=\!0.05/9\!\approx\!0.0056$)
none of the row-best margins anywhere in the appendix is significant.}
\label{tab:lpe_molpcba_k16}
\begin{tabular}{rcccc}
\toprule
$h$  & \texttt{map}              & \texttt{oap}              & \texttt{random\_augmented} & Welch $p$ (best vs.\ 2nd) \\
\midrule
 16  & $\mathbf{0.0716 \pm 0.0067}$ & $0.0703 \pm 0.0070$       & $0.0691 \pm 0.0008$         & $\approx 0.83$ \\
128  & $\mathbf{0.1923 \pm 0.0005}$ & $0.1836 \pm 0.0025$       & $0.1916 \pm 0.0009$         & $\approx 0.32$ \\
512  & $0.2468 \pm 0.0021$         & $0.2419 \pm 0.0028$       & $\mathbf{0.2483 \pm 0.0026}$  & $\approx 0.48$ \\
\bottomrule
\end{tabular}
\end{table}

\paragraph{Test-time averaging $\Delta$ is small.}
Across all nine $(k,h)$ cells the difference between the single-pass
deterministic-eigsh evaluation of the \texttt{random\_augmented} model
and its \textsc{aug-$K$} averaged evaluation
(Table~\ref{tab:lpe_molpcba_main}, right-most two columns) lies in
$[-0.0002, +0.0005]$~AP, within the seed noise at every cell. The
\textsc{aug-$K$} $\Delta$ is an aggregate metric, not a direct
measurement of model invariance; we test invariance directly in the
next paragraph.

\paragraph{Per-graph orbit stability.}
For each test graph we draw $K\!=\!64$ i.i.d.\ elements $g_j$ of the
sign-flip subgroup of $\G_{\mathrm{LapPE}}$, forward-pass the trained
\texttt{random\_augmented} model on each $g_j(\mathrm{LapPE})$, and
record the per-task predicted probability. The resulting per-graph
spread (max $-$ min predicted probability across the $K\!=\!64$ draws)
is a direct measure of how much the model's output varies across
orbit elements. Aggregating across all $43{,}793$ \texttt{ogbg-molpcba}
test graphs and $128$ tasks, at hidden dim $h\!=\!128$ we observe a
clean monotone scaling with $k$
(Table~\ref{tab:lpe_molpcba_orbit}): the per-graph spread is
$0.0018$ at $k\!=\!3$, $0.0034$ at $k\!=\!8$, and
$0.0048\!\pm\!0.0008$ at $k\!=\!16$ (mean over $3$ seeds at $k\!=\!16$;
single seed at $k\!\in\!\{3,8\}$). This is qualitatively consistent
with $|\G_{\mathrm{LapPE}}|\!=\!2^{k}$: more orbit elements give more
room to vary, and the model is closer to invariant for smaller $k$.
Concretely at $k\!=\!16$/$h\!=\!128$, the per-draw test AP across
$K\!=\!64$ orbit elements has standard deviation $0.0006$ within a
seed and total range $\approx\!0.0024$, comparable to the
seed-to-seed std of single-pass evaluation in
Table~\ref{tab:lpe_molpcba_k16} ($0.0009$). The Reynolds-style
(mean-probability) AP across the $K\!=\!64$ draws is
$0.1923\!\pm\!0.0017$ over seeds, within seed noise of the
matched-compute single-pass number ($0.1916\!\pm\!0.0009$ for
\texttt{random\_augmented} and $0.1923\!\pm\!0.0005$ for the
matched-compute \texttt{map} arm at the same cell). We read this as
direct evidence that the \texttt{random\_augmented} model is
approximately $\G_{\mathrm{LapPE}}$-invariant on the test distribution
at the configurations tested.

\begin{table}[h]
\centering
\caption{Per-graph orbit-stability of the trained
\texttt{random\_augmented} model under $K\!=\!64$ uniform i.i.d.\
draws from the sign-flip subgroup of $\G_{\mathrm{LapPE}}$, evaluated
on the \texttt{ogbg-molpcba} test set. ``Per-graph spread'' is the
mean over graphs and tasks of $(\max - \min)$ predicted probability
across the $K\!=\!64$ orbit elements; ``per-draw AP std (within
seed)'' is the standard deviation of the OGB AP scored on each single
orbit draw. The $k\!=\!16$ row is the mean $\pm$ sample std over the
$3$ training seeds; the $k\!\in\!\{3,8\}$ rows are single-seed
estimates ($s\!=\!0$). Per-graph spread scales monotonically with
$|\G_{\mathrm{LapPE}}|\!=\!2^{k}$.}
\label{tab:lpe_molpcba_orbit}
\begin{tabular}{rrcccc}
\toprule
$k$ & $h$ & seeds & per-graph spread & per-draw AP std (within seed) & Reynolds AP \\
\midrule
 3  & 128 & 1 & $0.0018$               & $0.0004$  & $0.1741$ \\
 8  & 128 & 1 & $0.0034$               & $0.0005$  & $0.1806$ \\
16  & 128 & 3 & $0.0048 \pm 0.0008$    & $0.0006$  & $0.1923 \pm 0.0017$ \\
\bottomrule
\end{tabular}
\end{table}

\paragraph{Reading the tables.}
Reading along the hidden-dim axis at fixed $k$ in
Table~\ref{tab:lpe_molpcba_main}, the row-best AP grows monotonically
with $h$ at every $k$, and the row-best arm is most often
\texttt{random\_augmented} (six of nine cells), with \texttt{map}
nominally row-best at $k\!=\!16$/$h\!\in\!\{16,128\}$ and \texttt{oap}
nominally row-best at $k\!=\!3$/$h\!=\!512$. The matched-compute
$k\!=\!16$ slice in Table~\ref{tab:lpe_molpcba_k16} resolves these
margins against seed variance: the smallest sample std on each row
is consistently the \texttt{random\_augmented} std ($0.0008$,
$0.0009$, $0.0026$), while \texttt{map} carries the largest std at
small $h$ ($0.0067$ at $h\!=\!16$). Under a two-sided Welch
$t$-test for unequal variances, the row-best versus row-second-best
margins at $k\!=\!16$ have $p$-values of $0.83$, $0.32$, and $0.48$
at $h\!=\!16,128,512$ respectively, none significant; this is the
honest reading of the appendix at $n\!=\!3$ seeds. At $k\!\in\!\{3,8\}$
the row-best margins are similar in magnitude or smaller, with
\texttt{random\_augmented} winning five of six rows there but at
margins of the same order as the seed std. We therefore present
Table~\ref{tab:lpe_molpcba_main} as a nominal-rank summary, not as
a benchmark claim.

\paragraph{Tie to the theory.}
On \texttt{ogbg-molpcba} at matched compute, with $3$ seeds per cell,
we did not detect a statistically significant difference between
training-time random sign-flip augmentation and projection-based
deterministic canonization at any of the nine $(k,h)$ cells; this is
qualitatively consistent with~\eqref{eq:coveringIneq} of
Proposition~\ref{propGen}, which suggests augmentation can match
canonization on tasks where the covering number of $\G$ is small,
but it is not a positive test of the inequality, since AP gaps of
order $10^{-3}$ are dominated by optimisation and seed noise rather
than by the input-space covering complexity the inequality bounds.
The \textsc{aug-$K$} column moves AP by at most $+0.0005$ across the
nine cells; this is consistent with the cached deterministic
eigsh-sign PE already lying close to the orbit-mean the network was
trained on, but is not itself a direct test of model invariance.

\section{Experimental Setup and Hyperparameter Details}
\label{app:exp_setup}

\subsection{Experimental Setup for the covering number experiment}
We applied the following preprocessing for \emph{each} sample: we first sample $256$ points, then we shift it to the positive axis, and then divide by the maximum axis along $x,y,z$, obtaining all entries to be positive and between $0$ and $1$. The metric we considered on the point cloud space is 
$$\rho_{\G}(X,Y) = \frac{1}{n}\sum^{n}_{i=1}\|X_i-Y_i\|_2 .$$
The same metric is used for the two canonizations, and its quotient $\rho_\G$ is used for the quotient space. 
%  so we took another approach. For each test sample $s$ we compute $d_s := min_{t \in Train}\rho(s,t)$, where Train are all training samples with same label as $s$, and then we average i.e. $\frac{1}{|\D_{test}}\sum_{s\in D_{test}}d_s$. We consider four types of distances: 
% \begin{enumerate}
%     \item $\rho_{reg}(X,Y) = \frac{1}{n}\sum^{n}_{i=1}|X_i-Y_i|$
%     \item $\rho_{wass}(X,Y)=min_{\pi \in S_n}\frac{1}{n}\sum^{n}_{i=1}|X_i-Y_{\pi(i)}|$
%     \item $\rho_{sort}(X,Y) = \frac{1}{n}\sum^{n}_{i=1}|sort(X)_i-sort(Y)_i|$
%     \item $\rho_{Hilbert}(X,Y) = \frac{1}{n}\sum^{n}_{i=1}|Hilbert(X)_i-Hilbert(Y)_i|$
% \end{enumerate}

\subsection{Experiment setup: ModelNet40 classification}
\begin{table}[ht]
\centering
\caption{Test accuracy of the Global MLP under varying training set sizes. The performance gap between Hilbert and lexicographical sorting widens as data becomes scarcer, highlighting the regularizing effect of optimal canonization.}
\label{tab:data_scarcity}
\begin{tabular}{lccc}
\toprule
\textbf{Training Samples} & \textbf{None} $\uparrow$ & \textbf{Lexicographical} $\uparrow$ & \textbf{Hilbert} $\uparrow$ \\
\midrule
9840 (full) & 0.433 & 0.700 & \textbf{0.748} \\
4920 & 0.408 & 0.641 & \textbf{0.714} \\
2460 & 0.343 & 0.573 & \textbf{0.673} \\
1230 & 0.307 & 0.504 & \textbf{0.601} \\
\bottomrule
\end{tabular}
\end{table}
In the ModelNet40 classification tasks, we sampled from each input CAD model point cloud, which consists of $1024$ points, and applied an initial batched affine normalization to $[0, 1]^3$. The data is then ordered using one of three strategies: no canonization (baseline), lexicographical sorting (i.e., sorting points sequentially by their $x$, then $y$, then $z$ coordinates), or Hilbert curve sorting (using a space size of $2^{12}$). Following ordering, the coordinates are projected into a higher-dimensional space using Random Fourier Features to capture high-frequency geometric details. The feature vector is then passed through a bottlenecked MLP consisting of residual blocks with dimensions $[256, 128, 64]$, yielding the final classification logits. The MLPs were trained for 100 epochs with a fixed batch size of 256, and were trained using AdamW.

To ensure a rigorous and fair comparison between the different ordering strategies, we conducted comprehensive grid sweeps to identify the optimal hyperparameters for the uncanonized baseline (\texttt{ply}), lexicographical sorting (\texttt{lex}), and Hilbert curve sorting (\texttt{hilbert}). The specific optimal hyperparameters extracted from our sweeps and utilized for the final evaluation are detailed in Table~\ref{tab:hyperparams}.
The best hyperparameters for each model on ModlenNet40 are used for training the models on ModelNet10. 
\begin{table}[h]
\centering
\caption{Optimal hyperparameters identified via grid search for each ordering strategy.}
\label{tab:hyperparams}
\begin{tabular}{lcccccc}
\toprule
\textbf{Method} & \textbf{Learning Rate} & \textbf{Dropout} & \textbf{Weight Decay} & \textbf{Fourier Scale} & \textbf{Label Smoothing} & \textbf{Bands} \\
\midrule
No canon. & $1.63 \times 10^{-3}$ & 0.3 & 0.05 & 1.0 & 0.1 & 2 \\
Lexsort  & $6.41 \times 10^{-4}$ & 0.1 & 0.05 & 0.1 & 0.2 & 3 \\
Hilbert  & $9.61 \times 10^{-4}$ & 0.3 & 0.01 & 0.1 & 0.2 & 3 \\
\bottomrule
\end{tabular}
\end{table}

\paragraph{Robustness Across Seeds}
To account for variance in weight initialization and optimization dynamics, we evaluated the final optimal configurations across 5 independent random seeds. The averaged test accuracy and standard deviations are reported in Table~\ref{tab:seed_robustness}. The results demonstrate that the Hilbert curve canonization not only achieves the highest average accuracy but also exhibits the lowest variance across different initializations.

\begin{table}[h]
\centering
\caption{Test accuracy across 5 independent random seeds using the optimal hyperparameters. Hilbert canonization yields both higher performance and greater stability.}
\label{tab:seed_robustness}
\begin{tabular}{lcc}
\toprule
\textbf{Method} & \textbf{Average Accuracy} $\uparrow$ & \textbf{Standard Deviation} \\
\midrule
No canon. & 0.433 & 0.0249 \\
Lexsort  & 0.700 & 0.01358 \\
Hilbert  & \textbf{0.748} & 0.00944 \\
\bottomrule
\end{tabular}
\end{table}

%\paragraph{Experiemental Setup: Modelnet10 experiment}
%\nd{yonatan: can you campare vs. Modelnet40 description above and think what to write here?}
%In the second set of experiments, we compare the performance of several invariant models on the official ModelNet10 train/test split. 
%From the original training split, we create a validation set using a fixed 80/20 random split, while leaving the test split unchanged.

%Each sample is a point cloud $X \in \mathbb{R}^{P \times 3}$ with $P=1024$ points.
%For every sample, we precompute and store three permutations of $\{1,\dots,P\}$: the identity permutation (\texttt{id}), coordinate-based sorting (\texttt{sort}), and Hilbert-based sorting (\texttt{hilbert}).
%During training and evaluation, we choose one permutation mode and apply it to the input cloud before feeding it to the model.

%The reordered cloud is then flattened into a vector in $\mathbb{R}^{3P}$ and classified by an MLP trained with cross-entropy loss over 10 classes.
%We optimize all models with Adam for 100 epochs using batch size 64.
%For each hyperparameter configuration, we save the checkpoint with the highest validation accuracy (\texttt{val\_acc}) and report test performance using that best-validation checkpoint.


\subsection{Experimental Setup: Modelnet 40 Data scarcity experiment}
To simulate data scarcity, we subsample the ModelNet40 training set using strides of 1, 2, 4, and 8, resulting in training subsets of 9840, 4920, 2460, and 1230 samples, respectively. Results are shown in Table  \ref{tab:data_scarcity}.

\subsection{Experimental Setup: PCA group averaging}
\label{subsec:ablation_pca}
We use a DeepSets architecture that is fully invariant to point order permutations but strictly non-invariant to the $\{-1, 1\}^3$ sign flips. We intentionally employ a lightweight, low-capacity version of this architecture. A highly over-parameterized model might achieve higher absolute accuracies, potentially masking the generalization gaps we aim to observe; by constraining the network capacity, we ensure that any performance differences are strictly attributable to the chosen alignment strategy. Each input point cloud consists of $1024$ points from the ModelNet40 dataset. We compare four approaches: (1) \textbf{Pure PCA}, which leaves the sign ambiguity unresolved (acting as a baseline or potentially \emph{poor} canonization due to eigensolver discontinuities); (2) \textbf{Skewness}, a deterministic canonization mapping $c: K \to K$ that resolves the ambiguity by orienting axes based on the third moment; (3) \textbf{Frame Averaging}, an explicitly $\G$-invariant model constructed by pooling the network's predictions across all $8$ elements of the group using log-sum-exp; and (4) \textbf{Random Frame}, which applies elements of $\G$ as random data augmentation during training. All models are optimized using Adam for 1600 epochs and evaluated across 5 different random seeds.

\begin{figure}[htbp]
    \centering
    \includegraphics[width=\linewidth]{figs/loss_pca.png}
    \caption{\textbf{Rotation Symmetry Accuracy.} The curves illustrate the comparative performance of the four approaches: Random Frame (blue), Frame Averaging (red), Skewness (orange), and the Pure PCA baseline (green). Frame Averaging achieves the strongest overall performance. While Random Frame and Skewness reach comparable final results—both significantly outperforming the Pure PCA baseline—the augmentation-based Random Frame method exhibits a noticeably slower convergence rate than deterministic canonization via Skewness.}
    \label{fig:loss_pca}
\end{figure}


\subsection{Experimental Setup: Rotated MNIST}
In the last set of the experiments, we compare the performance of several image rotation invariant models. We considered three models: simple CNN, CN(p4)-CNN, and AvgCNN. The backbone for all three models is the same and additional layers are added for CN(p4)-CNN, AvgCNN. 
In order to have the same setting, including the data split we rerun all models, including CN(p4)-CNN.
We used learning rate of $1e^{-3}$ and a batch size of $256$ and an AdamW optimizer. Each experiment is run for $5$ different seeds, and mean and standard deviation are reported. 
\newpage
\input{checklist.tex}
\end{document}


